\documentclass[11pt]{article}
\usepackage{tcee}
\addbibresource{3B39.bib}
\tcetema[Mercado único y competencia en la UE]{3.B.39}{El mercado único de la Unión Europea. El principio de libre circulación de mercancías, servicios, personas y capitales. La política de competencia.}
\tcefecha{19/11/2026}
\begin{document}
\tceetitulo

\section*{Introducción}
\minutos{2,5}

\begin{itemize}
\item \textbf{\emph{Enganche}}:
  \begin{itemize}
  \item «El pez grande se come al chico», dice el refrán. En 1988, el informe \autor{Cecchini} calculó que en más de la mitad de las ramas de la industria europea cabían 20~empresas de tamaño eficiente en un mercado comunitario, frente a solo 4~en el mayor de los mercados nacionales: los peces europeos eran pequeños porque nadaban en estanques pequeños\nocite{Cecchini1988Comision}.
  \item 36~años después, el informe \autor{Draghi} (2024) volvió al refrán desde el otro lado: frente a las empresas de Estados Unidos y de China, las europeas siguen siendo pequeñas, y pidió que la política de competencia mirara más la escala y la innovación\nocite{Draghi2024Future}. En 2026, la Comisión prepara unas nuevas directrices de concentraciones que lo recogen\nocite{CE2026DirectricesBorrador}.
  \item La pregunta que recorre la exposición es: \emph{¿cómo agranda el mercado único el estanque para que los peces crezcan, y cómo impide la política de competencia que el pez grande se coma al chico?}
  \end{itemize}

\item \textbf{\emph{Relevancia}}:
  \begin{itemize}
  \item \emph{Es el núcleo económico de la Unión}: 27~Estados, unos 450~millones de consumidores, 26~millones de empresas y un PIB de 18~billones de euros, al que se suman Islandia, Liechtenstein y Noruega por el Espacio Económico Europeo. La Comisión estima que eleva el PIB de la Unión al menos un 3-4~\% y que ha creado 3,6~millones de empleos\nocite{CE2025EstrategiaFicha}.
  \item \emph{Está lejos de completarse}: el Fondo Monetario Internacional (FMI) estima que las barreras que quedan dentro de Europa equivalen a un arancel del 44~\% en las manufacturas y de hasta el 110~\% en los servicios\nocite{FMI2024REO}.
  \item \emph{Para España}: más de la mitad de sus exportaciones de bienes van a otros Estados de la Unión, y su política de competencia, sus ayudas públicas y buena parte de su regulación se deciden o se controlan en Bruselas \vertema{4.A.15} \vertema{4.A.22}.
  \end{itemize}

\item \textbf{\emph{Contextualización}}:
  \begin{itemize}
  \item \emph{Qué es}: el \concepto{mercado interior} es «un espacio sin fronteras interiores, en el que la libre circulación de mercancías, personas, servicios y capitales estará garantizada» (art. 26~del Tratado de Funcionamiento de la Unión Europea, TFUE)\nocite{UE2016TFUE}. Es una fase de la integración más profunda que la unión aduanera y menos que la unión económica \vertema{3.B.10}.
  \item \emph{2 formas de integrar}, en la terminología de \autor{Jan Tinbergen} (1954)\nocite{Tinbergen1954International}:\footnote{\autor{Tinbergen} recibió el primer Premio Nobel de Economía en 1969, con \autor{Ragnar Frisch}, «por haber desarrollado y aplicado modelos dinámicos para el análisis de los procesos económicos»\nocite{Nobel1969FrischTinbergen}.}
    \begin{itemize}
    \item la \concepto{integración negativa}: quitar barreras (aranceles, controles en frontera, normas que discriminan);
    \item la \concepto{integración positiva}: crear reglas e instituciones comunes (armonización, política de competencia) para que el mercado sin barreras funcione.
    \end{itemize}
  \item \emph{Cómo se reparten las competencias}: el mercado interior es una competencia \emph{compartida} entre la Unión y los Estados (art. 4~TFUE); las normas sobre competencia necesarias para su funcionamiento son una de las 5~competencias \emph{exclusivas} de la Unión (art. 3~TFUE) \vertema{3.B.36}.
  \end{itemize}

\item \textbf{\emph{Problemática}}:
  \begin{itemize}
  \item ¿Por qué conviene un mercado único, qué dicen los Tratados y cómo se ha construido?
  \item ¿Qué garantiza cada una de las 4~libertades y qué barreras quedan?
  \item ¿Por qué hace falta una política de competencia común y cómo controla a las empresas y a los Estados?
  \item ¿Debe la Unión aceptar empresas más grandes para competir en el mundo?
  \end{itemize}

\item \textbf{\emph{Estructura}}:
  \begin{itemize}
  \item La exposición va del mercado en su conjunto a cada una de sus libertades y de ahí a las reglas que lo protegen:
    \begin{itemize}
    \item \concepto{I. El mercado único} (\ref{sec:mercado-unico}): por qué conviene (la competencia y la escala de la integración de mercados oligopolísticos), qué dicen los Tratados (el art. 26~TFUE y el 114) y cómo se ha construido, de Roma a la estrategia de 2025.
    \item \concepto{II. Las 4~libertades} (\ref{sec:libertades}): mercancías (con las barreras fiscales), servicios, personas y capitales, con su justificación, su base jurídica, su jurisprudencia y lo que falta.
    \item \concepto{III. La política de competencia} (\ref{sec:competencia}): su justificación, su marco jurídico y sus 4~pilares (acuerdos, abusos, concentraciones y ayudas de Estado), con los instrumentos nuevos y las ideas que los justifican: el coste de los cárteles y la clemencia, el compromiso de \autor{Williamson} y la carrera de subvenciones.
    \end{itemize}
  \end{itemize}
\end{itemize}

\begin{esquema}[unbreakable]
\textbf{3.B.39. El mercado único de la UE y la política de competencia}\\[2pt]
Introducción\esqtiempo{2,5}\\
I. El mercado único: justificación, marco jurídico y evolución\esqtiempo{7}\\
\hspace*{1.5em}1. La justificación económica: competencia y escala\\
\hspace*{1.5em}2. Fundamentos jurídicos básicos del mercado interior\\
\hspace*{1.5em}3. De Roma a la estrategia de 2025\\
II. Las 4~libertades de circulación\esqtiempo{9}\\
\hspace*{1.5em}1. Mercado interior de bienes\\
\hspace*{1.5em}2. Mercado interior de servicios\\
\hspace*{1.5em}3. Libre circulación de personas\\
\hspace*{1.5em}4. Libre circulación de capitales\\
III. La política europea de competencia\esqtiempo{10}\\
\hspace*{1.5em}1. Motivación de la teoría económica\\
\hspace*{1.5em}2. Marco jurídico: los 4~pilares\\
\hspace*{1.5em}3. Promoción, coordinación y debate actual\\
Conclusión\esqtiempo{1,5}
\end{esquema}

\begin{anotaciones}
\textbf{Mapa de la exposición.} Cada bloque responde a una pregunta y deja abierta la del siguiente; a la derecha, lo que se cuenta en cada subapartado.
\begin{infografiacaja}
\caption{Mapa del tema: la pregunta de cada bloque y el hilo que los une}\label{figu:mapa-tema}
\grafico[\linewidth]{mapa-tema}
\fuente{elaboración propia.}
\end{infografiacaja}
\end{anotaciones}


\begin{anotaciones}
\textbf{Notación.} En el modelo del mercado integrado, $p(Q)=a-bQ$ es la demanda inversa de cada mercado, con $b=1/S$ y $S$ el \emph{tamaño del mercado} (símbolo local de este tema); $c$ y $\mathrm{CF}$ son el coste marginal constante y el coste fijo de cada empresa, $\mathrm{CMe}(q)=c+\mathrm{CF}/q$; $n$ es el número de empresas y $q_j$ la producción de la empresa $j$. En el coste del cártel, $\Delta B$ es el beneficio ilícito, $\pi$ la probabilidad de que se descubra y se sancione el cártel y $M$ la multa; en su sostenibilidad, $\beta$ es el factor de descuento y $B^{\mathrm C}$, $B^{\mathrm D}$ y $B^{\mathrm N}$ los beneficios por periodo de coludir, de desviarse y del castigo, en redonda, como en la teoría del oligopolio \vertema{3.A.19}. En el compromiso de \autor{Williamson}, $p_0=\mathrm{CMe}_0$ y $p_1$ son los precios antes y después de la fusión, $\mathrm{CMe}_1$ el coste medio después y $\varepsilon$ la elasticidad de la demanda. TFUE: Tratado de Funcionamiento de la Unión Europea.
\end{anotaciones}

% ==========================================================================
\section{El mercado único:\cortetitulo{} justificación, marco jurídico y evolución}\label{sec:mercado-unico}
\minutos{7}

\begin{itemize}
\item Empezamos por el mercado único en su conjunto: por qué conviene a la Unión, qué dicen los Tratados y cómo se ha construido.
\item \emph{Conclusión anticipada}: su principal ganancia no es la clásica de la especialización, sino la de la competencia y la escala: con mercados más grandes hay empresas más grandes y márgenes más pequeños. Los Tratados lo definen en el art. 26~TFUE y lo construyen sobre todo con el art. 114; se hizo en 2~tiempos, la unión aduanera de 1968 y el programa de 1992, y 3~décadas después sigue incompleto, sobre todo en los servicios.
\end{itemize}

% --------------------------------------------------------------------------
\subsection{La justificación económica: competencia y escala}\label{sec:efectos}

\begin{itemize}
\item \emph{Los argumentos de la teoría}: la liberalización del comercio y de los factores dentro de un mercado común se apoya en varias teorías, que se desarrollan en sus temas:
  \begin{itemize}
  \item \emph{la ventaja comparativa}: la producción se hace donde es más eficiente y crecen las posibilidades de consumo (\autor{Haberler}, 1936) \vertema{3.B.5};
  \item \emph{la nueva teoría del comercio}: economías de escala y variedad (\autor{Krugman}, 1979) \vertema{3.B.6};
  \item \emph{la novísima}: la apertura reasigna la producción hacia las empresas más productivas (\autor{Melitz}, 2003) \vertema{3.B.6};
  \item \emph{la teoría de la integración}: creación y desviación de comercio (\autor{Viner}, 1950) y efectos dinámicos \vertema{3.B.10};
  \item \emph{la geografía económica}: la integración puede concentrar la actividad en un centro y vaciar la periferia, de ahí la política de cohesión \vertema{3.B.10} \vertema{3.B.40};
  \item \emph{el crecimiento endógeno}: más variedades de bienes intermedios y más difusión del conocimiento (\autor{Romer}, 1990) \vertema{3.B.9}.
  \end{itemize}
\item Entre economías parecidas, como las europeas, la ganancia más importante no es la especialización, sino la que el informe \autor{Cecchini} (1988) puso en el centro: \emph{más competencia y más escala}.
\end{itemize}

\subsubsection{La integración de mercados oligopolísticos}\label{sec:modelo-integracion}

\begin{itemize}
\item \emph{La idea}: con mercados nacionales separados, cada uno es pequeño: caben pocas empresas, cada una produce poco (lejos de su escala eficiente) y, como compite con pocas rivales, cobra un margen alto. Al integrarse los mercados ocurren 2~cosas:
  \begin{itemize}
  \item cada empresa compite con más rivales y percibe una demanda más elástica, y baja el margen: es el \concepto{efecto procompetitivo};
  \item con márgenes más bajos no caben todas las empresas: salen algunas, y las que quedan producen más y más barato: es el \concepto{efecto de racionalización}.
  \end{itemize}
\item \emph{El resultado}: es la lógica de \textcite{SmithVenables1988Completing}, que simularon el mercado único con oligopolios de \autor{Cournot}, y de \textcite{Krugman1979Increasing} con competencia monopolística.\footnote{\autor{Paul Krugman} recibió el Premio Nobel de Economía en 2008 «por su análisis de los patrones del comercio y de la localización de la actividad económica»\nocite{Nobel2008Krugman}.} En su versión más sencilla, con empresas iguales y libre entrada (anexo~\ref{anexo:integracion}), el margen sobre el coste marginal de equilibrio depende del coste fijo y del tamaño del mercado:
  \funcionclave[eq:margen-tamano]{p^{*}-c=\sqrt{\frac{\mathrm{CF}}{S}}}
  El margen baja con la raíz cuadrada del tamaño del mercado: si 2~mercados iguales se integran, cae un 29~\%, cada empresa produce un 41~\% más y quedan menos empresas que sumando las de los 2~países, aunque cada consumidor puede comprar a más que antes.
\item \emph{Las implicaciones}:
  \begin{itemize}
  \item \emph{la ganancia es de eficiencia productiva y asignativa a la vez}: precios más bajos para los consumidores y costes medios más bajos por la escala; por eso \autor{Cecchini} daba ganancias varias veces mayores que las de la unión aduanera;
  \item \emph{la integración trae concentración}: menos empresas y más grandes; no es un fallo, sino el mecanismo por el que se aprovecha la escala, pero exige una política de competencia que impida que las empresas que quedan se pongan de acuerdo o se fusionen hasta recuperar el margen (bloque~\ref{sec:competencia});
  \item \emph{las barreras que no son aranceles también segmentan}: normas técnicas distintas, licencias nacionales o compras públicas cerradas mantienen el mercado de cada empresa en el tamaño de su país, y quitarlas tiene el mismo efecto que integrar.
  \end{itemize}
\item \emph{Las extensiones}: con comercio cruzado de un mismo bien (el «\emph{dumping} recíproco» de \autor{Brander} y \autor{Krugman}, 1983), los márgenes bajan ya con mercados segmentados, pero las mayores ganancias llegan cuando las empresas dejan de tratar cada país como un mercado aparte \vertema{3.B.10}; con productos diferenciados, el mercado más grande ofrece además más variedades, y con empresas heterogéneas expulsa a las menos productivas \vertema{3.B.6}.
\end{itemize}

\subsubsection{La evidencia}\label{sec:efectos-evidencia}

\begin{itemize}
\item \emph{Comercio}: el comercio de bienes dentro de la Unión pasó del 17~\% del PIB en 2002 al 22~\% en 2025, con un pico del 26~\% en 2022 por los precios de la energía; el de servicios, del 4,8~al 8~\% entre 2010 y 2025 (gráfico~\ref{fig:comercio-intra}). En el último informe anual del mercado único, el comercio intra-UE en porcentaje del PIB es uno de los indicadores que empeoran\nocite{CE2026InformeAnual}.
\item \emph{Márgenes}: \textcite{Badinger2007Has} estima que los márgenes sobre el coste marginal bajaron tras el programa de 1992 en la industria y la construcción, pero no en buena parte de los servicios, donde el mercado siguió fragmentado.
\item \emph{PIB}: las estimaciones a posteriori van desde el 3-4~\% del PIB de la Comisión\nocite{CE2025EstrategiaFicha} hasta en torno al 9~\% de los modelos que suman bienes y servicios\nocite{inTVeld2019Economic}; \textcite{MayerVicardZignago2019Cost} estiman que el mercado único ha más que duplicado el comercio de bienes entre sus miembros y que reducirlo a un acuerdo comercial convencional costaría a cada Estado, de media, un 4,4~\% de bienestar (entre el 3~y el 7~\% según el escenario).
\item \emph{Lo que queda}: el FMI estima que en 2020 las barreras dentro de Europa equivalían a un arancel del 44~\% en la manufactura media (frente al 15~\% entre los estados de Estados Unidos) y de hasta el 110~\% en los servicios, y que llevarlas al nivel estadounidense elevaría la productividad casi un 7~\% a largo plazo\nocite{FMI2024REO,FMI2024Kammer}.
\item \emph{Valoración}: el argumento explica por qué la ganancia del mercado único se mide en competencia y escala, pero supone que el mercado integrado sigue siendo competitivo: no dice nada de cárteles, de abusos, de fusiones que eliminan rivales ni de Estados que subvencionan a sus empresas. Esas son las tareas de la política de competencia (bloque~\ref{sec:competencia}).
\end{itemize}

\begin{figure}[H]
\centering
\caption{Comercio dentro de la Unión Europea, 2002-2025}\label{fig:comercio-intra}
\grafico[0.88\textwidth]{comercio-intra-ue}
\fuente{Eurostat (tablas \texttt{ext\_lt\_intratrd}, \texttt{bop\_its6\_det} y \texttt{nama\_10\_gdp}; consulta del 14/11/2026)\nocite{Eurostat2026Comercio}. Bienes: media de exportaciones e importaciones dentro de la UE-27 en porcentaje del PIB; servicios: media de créditos y débitos con la UE-27 (2025, provisional).}
\end{figure}

% --------------------------------------------------------------------------
\subsection{Fundamentos jurídicos básicos del mercado interior}\label{sec:fundamentos}

\begin{itemize}
\item Una vez visto por qué conviene un mercado único, pasamos a lo que dicen los Tratados: quién es competente, sobre qué base se actúa y cómo se decide.
\item \emph{La base jurídica: el art. 26~TFUE}. Es el artículo que define el mercado interior y da a la Unión el mandato de construirlo. Procede del art. 8A que introdujo el Acta Única en 1986, con el plazo del 31~de diciembre de 1992, y dice\nocite{UE2016TFUE}:
\end{itemize}

\begin{ideaclave}
\textbf{Artículo 26~del Tratado de Funcionamiento de la Unión Europea}
\begin{enumerate}
\item La Unión adoptará las medidas destinadas a establecer el mercado interior o a garantizar su funcionamiento, de conformidad con las disposiciones pertinentes de los Tratados.
\item El mercado interior implicará un espacio sin fronteras interiores, en el que la libre circulación de mercancías, personas, servicios y capitales estará garantizada de acuerdo con las disposiciones de los Tratados.
\item El Consejo, a propuesta de la Comisión, definirá las orientaciones y condiciones necesarias para asegurar un progreso equilibrado en el conjunto de los sectores considerados.
\end{enumerate}
\end{ideaclave}

\begin{itemize}
\item \emph{Lo que dice el art. 26}, en 3~ideas:
  \begin{itemize}
  \item \emph{un mandato} (26.1): la Unión no solo puede, sino que debe adoptar las medidas que establezcan el mercado interior y lo hagan funcionar;
  \item \emph{una definición} (26.2): un espacio \emph{sin fronteras interiores} con las \emph{4 libertades}; cada libertad se garantiza «de acuerdo con las disposiciones de los Tratados», es decir, con su propio título y sus propias excepciones;
  \item \emph{un equilibrio} (26.3): el Consejo vela por que el mercado avance a la vez en todos los sectores.
  \end{itemize}
\item \emph{Del art. 26~cuelga todo lo demás} (figura~\ref{figu:base-juridica}), de más a menos peso:
  \begin{itemize}
  \item el \emph{art. 114}, la base general de la armonización, con la que se ha hecho casi todo el mercado interior;
  \item los \emph{títulos de cada libertad} (arts. 28-66), con sus excepciones (tabla~\ref{tab:libertades});
  \item las \emph{bases propias}, con unanimidad, de lo que el art. 114~excluye: la fiscalidad (arts. 113~y 115) \vertema{4.B.20}.
  \end{itemize}
\end{itemize}

\begin{infografia}[H]
\centering
\caption{La base jurídica del mercado interior: del art. 26~a cada libertad}\label{figu:base-juridica}
\grafico[\textwidth]{base-juridica}
\fuente{elaboración propia a partir del TFUE\nocite{UE2016TFUE}. CESE: Comité Económico y Social Europeo.}
\end{infografia}

\begin{itemize}
\item \emph{Delimitación de competencias}: el mercado interior es un ámbito de \concepto{competencia compartida} (art. 4.2.a TFUE). Según el art. 2.2~TFUE, la Unión y los Estados pueden legislar en él, y los Estados ejercen su competencia en la medida en que la Unión no haya ejercido la suya. Por eso rigen los principios de \emph{subsidiariedad} y \emph{proporcionalidad} (art. 5~del Tratado de la Unión Europea, TUE). Las normas sobre competencia necesarias para su funcionamiento son, en cambio, una competencia \emph{exclusiva} (art. 3.1.b TFUE), como la unión aduanera (art. 3.1.a) (anexo~\ref{anexo:competencias}) \vertema{3.B.36}.
\item \emph{La toma de decisiones: el art. 114~TFUE}. El Parlamento Europeo y el Consejo adoptan, por el \concepto{procedimiento legislativo ordinario} y previa consulta al Comité Económico y Social Europeo (CESE), las medidas de aproximación de las legislaciones que tengan por objeto el establecimiento y el funcionamiento del mercado interior:
  \begin{itemize}
  \item la Comisión propone, y el Parlamento y el Consejo, este por \emph{mayoría cualificada}, aprueban (art. 294~TFUE): es la novedad del Acta Única que desbloqueó el programa de 1992;
  \item la Comisión parte de un \emph{nivel de protección elevado} en salud, seguridad, medio ambiente y consumidores (art. 114.3);
  \item un Estado puede mantener o introducir normas más protectoras si las notifica y la Comisión las aprueba, porque no son una discriminación arbitraria ni una restricción encubierta (arts. 114.4~a 114.6).
  \end{itemize}
\item \emph{Lo que el art. 114~excluye} (art. 114.2): las disposiciones fiscales, las relativas a la libre circulación de personas y las de los derechos de los trabajadores, que tienen bases propias:
  \begin{itemize}
  \item la \emph{imposición indirecta} (art. 113) y la \emph{directa} (art. 115): el Consejo decide por \emph{unanimidad}, por un procedimiento legislativo especial y previa consulta al Parlamento y al CESE (\ref{sec:barreras-fiscales});
  \item las \emph{personas}: los arts. 21~y 46~a 48, con sus propios procedimientos (\ref{sec:personas}).
  \end{itemize}
\item \emph{El CESE}, órgano consultivo de la Unión, representa a empresarios, trabajadores y otras organizaciones de la sociedad civil; su dictamen es preceptivo, pero no vinculante.
\item \emph{Mercado común, interior y único}: el Tratado de Roma hablaba de \emph{mercado común}; el Acta Única introdujo \emph{mercado interior}, y \emph{mercado único} es el término político que se usa desde el programa de 1992. Hoy son sinónimos.
\end{itemize}

\begin{anotaciones}
\textbf{Las 4~libertades en el Tratado.} Cada libertad tiene su título en el TFUE, con una prohibición de las restricciones y unas excepciones tasadas que el Tribunal de Justicia interpreta de forma estricta (caben solo si son necesarias y proporcionadas):
\begin{itemize}
\item \emph{Mercancías}: unión aduanera, sin derechos de aduana ni exacciones de efecto equivalente y con un arancel común (arts. 28-32); prohibición de las restricciones cuantitativas y de las medidas de efecto equivalente (arts. 34-35); y prohibición de la discriminación fiscal (art. 110). Las excepciones son las del art. 36 (orden y seguridad públicos, salud, patrimonio, propiedad industrial) y las «exigencias imperativas» que añadió la jurisprudencia.
\item \emph{Personas}: el derecho de todo ciudadano a circular y residir (art. 21) y la libre circulación de los trabajadores sin discriminación por nacionalidad (arts. 45-48, con la coordinación de la seguridad social del art. 48). Las excepciones son el orden, la seguridad y la salud públicos (art. 45.3) y los empleos públicos con ejercicio de autoridad (art. 45.4).
\item \emph{Servicios y establecimiento}: la libertad de instalarse de forma estable en otro Estado (arts. 49-55) y la de prestar servicios desde el propio (arts. 56-62). Las excepciones son las actividades de poder público (art. 51) y el orden, la seguridad y la salud públicos (art. 52), que el art. 62~aplica también a los servicios; los transportes tienen su propio título.
\item \emph{Capitales y pagos}: la prohibición de toda restricción, también frente a terceros países (art. 63). Los Estados pueden distinguir por residencia en sus impuestos, supervisar y proteger el orden público (art. 65), y frente a terceros países caben las restricciones previas a 1994 (art. 64) y salvaguardias de 6~meses (art. 66).
\end{itemize}
\begin{tablacaja}
\centering
\caption{Las 4~libertades: base jurídica, excepciones y jurisprudencia básica}\label{tab:libertades}
\small
\begin{tabular}{>{\raggedright\arraybackslash}p{2.4cm}>{\raggedright\arraybackslash}p{3.2cm}>{\raggedright\arraybackslash}p{3.5cm}>{\raggedright\arraybackslash}p{4.3cm}}
\hline
\textbf{Libertad} & \textbf{Base jurídica (TFUE)} & \textbf{Excepciones} & \textbf{Sentencias clave}\\
\hline
Mercancías & arts. 28-37 (unión aduanera, 34-35 restricciones cuantitativas y medidas de efecto equivalente) y 110 & art. 36 y «exigencias imperativas» & \emph{Dassonville} (1974), \emph{Cassis de Dijon} (1979), \emph{Keck} (1993), \emph{fresas españolas} (1997)\\
Personas & art. 21 (ciudadanía), 45-48 (trabajadores) & art. 45.3 (orden, seguridad y salud públicos) y 45.4 (empleo público) & \emph{Bosman} (1995)\\
Servicios y establecimiento & arts. 49-55 (establecimiento) y 56-62 (servicios) & arts. 51 (poder público) y 52 (orden, seguridad y salud públicos) & \emph{Van Binsbergen} (1974), \emph{Laval} y \emph{Viking} (2007)\\
Capitales y pagos & arts. 63-66 & art. 65 (fiscalidad, supervisión, orden público); 64 y 66 (terceros países); 143-144 & \emph{acciones de oro}: Comisión contra España (2003)\\
\hline
\end{tabular}
\fuente{elaboración propia a partir del TFUE y de la jurisprudencia del Tribunal de Justicia\nocite{UE2016TFUE}.}
\end{tablacaja}
\end{anotaciones}

% --------------------------------------------------------------------------
\subsection{De Roma a la estrategia de 2025}\label{sec:historia}

\subsubsection{La unión aduanera y el estancamiento (1957-1985)}\label{sec:union-aduanera}

\begin{itemize}
\item \emph{El Tratado de Roma} (1957) creó la Comunidad Económica Europea con el objetivo de un \concepto{mercado común}. No lo definía, pero implicaba las 2~integraciones: la negativa (suprimir barreras) y la positiva (políticas comunes).
\item \emph{La unión aduanera} se completó el 1~de julio de 1968, 18~meses antes de lo previsto:
  \begin{itemize}
  \item se suprimieron los aranceles y las restricciones cuantitativas entre los 6, también para los productos de terceros países en libre práctica;\footnote{Un producto está en \emph{libre práctica} cuando ha cumplido las formalidades de importación y pagado los aranceles en un Estado miembro: desde entonces circula como uno comunitario.}
  \item se adoptó un \concepto{arancel exterior común}, calculado como la media aritmética de los aranceles de partida;\footnote{En realidad, de los de las 4~zonas arancelarias previas, porque Bélgica, los Países Bajos y Luxemburgo ya formaban el Benelux. Los países que se han adherido después adoptan el arancel común tras un periodo transitorio \vertema{3.B.41}.} la ganancia que se le atribuyó fue pequeña, de alrededor del 1~\% del PIB \vertema{3.B.10}.
  \end{itemize}
\item \emph{El estancamiento} (1973-1985): tras las crisis del petróleo, los Estados levantaron \concepto{barreras no arancelarias}: controles en frontera, normas técnicas distintas, compras públicas reservadas a las empresas nacionales. La regla de la unanimidad para armonizar bloqueaba cualquier avance. El Tribunal de Justicia fue quien mantuvo vivo el mercado común, con 2~sentencias:
  \begin{itemize}
  \item \emph{Dassonville} (1974) definió de forma muy amplia las «medidas de efecto equivalente» a una restricción cuantitativa, prohibidas por el Tratado;
  \item \emph{Cassis de Dijon} (1979) sentó el \concepto{principio de reconocimiento mutuo} (\ref{sec:bienes}).
  \end{itemize}
\end{itemize}

\subsubsection{El programa de 1992: Libro Blanco, Acta Única e informe Cecchini}\label{sec:programa-1992}

\begin{itemize}
\item \emph{El Libro Blanco} (1985)\footnote{Los \emph{libros blancos} de la Comisión contienen propuestas concretas de acción de la Unión en un ámbito; los \emph{libros verdes} solo abren la reflexión y la consulta sobre un tema.} de la Comisión presidida por \autor{Jacques Delors}, preparado por el comisario \autor{Lord Cockfield}, listó unas 280~medidas para eliminar 3~tipos de barreras antes del 31~de diciembre de 1992:
  \begin{itemize}
  \item \emph{físicas}: los controles en frontera de mercancías y personas;
  \item \emph{técnicas}: las normas y reglamentaciones nacionales distintas y las compras públicas cerradas;
  \item \emph{fiscales}: las diferencias de imposición indirecta, que obligaban a ajustes en frontera.
  \end{itemize}
\item \emph{El Acta Única Europea} (firmada en 1986, en vigor el 1~de julio de 1987) le dio la base jurídica: introdujo el concepto de mercado interior con el plazo de 1992 y, sobre todo, la \emph{mayoría cualificada} para la armonización (el actual art. 114~TFUE), salvo en fiscalidad.
\item \emph{El informe Cecchini} (1988): la Comisión encargó a un grupo de expertos dirigido por \autor{Paolo Cecchini} estimar los «costes de la no Europa»\nocite{Cecchini1988Comision,Emerson1988Economics,Muns2005Lecturas}:
  \begin{itemize}
  \item las barreras en frontera costaban en torno al 1,8~\% del valor del comercio intracomunitario, y las técnicas, algo menos del 2~\% de los costes de las empresas;
  \item a medio plazo, completar el mercado elevaría el PIB un \emph{4,5~\%}, bajaría los precios al consumo un 6,1~\%, crearía 1,8~millones de empleos, mejoraría el saldo público en 2,2~puntos del PIB y el exterior en 1~punto, con un margen de error del 30~\%\nocite{Cecchini1988Comision};
  \item la clave estaba en la \emph{competencia y la escala}: un tercio de la industria podría reducir costes entre el 1~y el 7~\%, y las economías de escala solas valían un 2~\% del PIB.
  \end{itemize}
\item \emph{El 1~de enero de 1993} nació el \concepto{mercado único}: desaparecieron los controles de mercancías en las fronteras interiores, con un IVA intracomunitario sin ajustes en frontera (\ref{sec:bienes}) y una estadística nueva del comercio (Intrastat).
\end{itemize}

\begin{infografia}[H]
\centering
\caption{El mercado único en el tiempo: de Roma a la estrategia de 2025}\label{figu:linea-tiempo}
\grafico[0.98\textwidth]{linea-tiempo}
\fuente{elaboración propia. IMERA: Instrumento de Emergencia y Resiliencia del Mercado Interior; UAI: Unión de Ahorros e Inversiones.}
\end{infografia}

\subsubsection{Los relanzamientos (2010-2026)}\label{sec:relanzamientos}

\begin{itemize}
\item \emph{Del informe \autor{Monti} a la estrategia de 2015}: tras la crisis financiera, el informe de \autor{Mario Monti} (2010) avisó de la «fatiga del mercado único»\nocite{Monti2010New}; siguieron 3~iniciativas:
  \begin{itemize}
  \item las Actas del Mercado Único, con 12~palancas en 2011 y 12~acciones clave en 2012\nocite{CE2011ActaMercadoUnico};
  \item la estrategia del mercado único digital (2015)\nocite{CE2015MercadoDigital};
  \item la Unión de Mercados de Capitales (2015)\nocite{CE2015UMC}.
  \end{itemize}
\item \emph{Los informes de 2024}:
  \begin{itemize}
  \item \autor{Enrico Letta}, «Mucho más que un mercado» (abril de 2024)\nocite{Letta2024Much}, propuso una \concepto{quinta libertad} (la investigación, la innovación, el conocimiento y los datos), una \concepto{Unión de Ahorros e Inversiones} que movilice los 33~billones de euros de ahorro privado, del que unos 300.000~millones al año salen hacia Estados Unidos, y un \concepto{28.º régimen}: un Derecho de sociedades europeo, optativo, junto a los 27~nacionales;
  \item \autor{Mario Draghi}, «El futuro de la competitividad europea» (9~de septiembre de 2024)\nocite{Draghi2024Future}, cifró la inversión adicional necesaria en 750.000-800.000 millones de euros al año (el 4,4-4,7~\% del PIB de 2023) y pidió que el control de concentraciones valorara el potencial de innovación futura (\ref{sec:concentraciones}).
  \end{itemize}
\item \emph{El Instrumento de Emergencia y Resiliencia del Mercado Interior} (IMERA, Reglamento (UE) 2024/2747), aplicable desde el 29~de mayo de 2026, responde a los cierres de fronteras de la pandemia: en una emergencia que active el Consejo, los Estados no pueden prohibir las exportaciones a otros Estados de bienes y servicios esenciales para la crisis, y la Comisión puede, como último recurso, pedir información a las empresas, dar prioridad a pedidos y comprar en común\nocite{UE2024IMERA}.
\item \emph{La estrategia del mercado único} (21~de mayo de 2025)\nocite{CE2025Estrategia,CE2025EstrategiaFicha} identificó las \concepto{10 barreras más dañinas} («\emph{Terrible Ten}»), el 60~\% de las cuales son las mismas que hace 20~años:
  \begin{itemize}
  \item normas de la Unión demasiado complejas;
  \item dificultades para crear y operar empresas;
  \item normalización demasiado lenta;
  \item falta de implicación de los Estados en aplicar las normas;
  \item reconocimiento limitado de las cualificaciones profesionales;
  \item reglas de envases, etiquetado y residuos fragmentadas;
  \item normas de producto anticuadas y productos que no las cumplen;
  \item trámites pesados para desplazar trabajadores;
  \item regulación nacional de los servicios restrictiva y divergente;
  \item restricciones territoriales de suministro, con las que los fabricantes impiden a los minoristas comprar en otro Estado.
  \end{itemize}
\item \emph{2026}: las 3~instituciones firmaron el 24~de abril la hoja de ruta «Una Europa, un mercado», que obliga a acordar antes de finales de 2027 las iniciativas pendientes (paquetes de simplificación «ómnibus», normalización, servicios, energía y digital)\nocite{CE2026UnaEuropa}. La Comisión propuso el 18~de marzo el 28.º régimen, \concepto{EU Inc.}, una sociedad optativa que se registra en línea en 48~horas, por menos de 100~euros y sin capital mínimo; el Consejo fijó su posición el 23~de octubre y el Parlamento aún no tiene la suya\nocite{CE2026EUInc}.
\end{itemize}

\begin{itemize}
\item \emph{Recapitulación del bloque~\ref{sec:mercado-unico}}: el mercado único se justifica por la competencia y la escala, que resume $p^{*}-c=\sqrt{\mathrm{CF}/S}$; el art. 26~TFUE lo define y obliga a construirlo, y el 114~lo armoniza por mayoría cualificada, salvo la fiscalidad; se hizo en la unión aduanera y en el programa de 1992, y se relanza desde 2024. Pasamos a ver cómo se ha conseguido en cada una de las 4~libertades.
\end{itemize}

% ==========================================================================
\section{Las 4~libertades de circulación}\label{sec:libertades}
\minutos{9}

\begin{itemize}
\item Una vez estudiado el mercado en su conjunto, pasamos a cada una de sus libertades, con el mismo esquema: justificación, base jurídica y jurisprudencia, derecho derivado e instrumentos, y balance.
\item \emph{Conclusión anticipada}: la libertad de mercancías es la más avanzada; la de capitales, la más amplia (vale también frente a terceros países); la de servicios, la más atrasada, y la de personas, la más limitada en la práctica por las lenguas, las cualificaciones y la protección social.
\end{itemize}

% --------------------------------------------------------------------------
\subsection{Mercado interior de bienes}\label{sec:bienes}

\subsubsection{Justificación y marco jurídico}\label{sec:bienes-marco}

\begin{itemize}
\item \emph{Justificación económica}: ¿por qué conviene que los bienes circulen sin barreras dentro de la Unión? La teoría del comercio da 4~razones, cada una desarrollada en su tema:
  \begin{itemize}
  \item \emph{las ganancias de la especialización}: cada país produce lo que hace relativamente mejor, por su productividad (\autor{David Ricardo}, 1817) o por su dotación de factores (\autor{Eli Heckscher} y \autor{Bertil Ohlin}), y puede consumir más de lo que produciría solo \vertema{3.B.5};
  \item \emph{las de la escala y la variedad}: entre países parecidos, como los europeos, domina el comercio intraindustrial, y la apertura da más variedades y empresas más grandes (\autor{Krugman}, 1979) y reasigna la producción hacia las más productivas (\autor{Melitz}, 2003) \vertema{3.B.6};
  \item \emph{las de quitar barreras que son puro coste}: un arancel crea una renta que recauda alguien, pero un control en frontera o una norma técnica distinta solo consumen recursos; por eso suprimir estas barreras «de fricción» da ganancias netas mayores que suprimir un arancel de la misma cuantía\nocite{BaldwinWyplosz2020Economics} \vertema{3.B.7};
  \item \emph{las de ir más allá de la unión aduanera}: una unión aduanera puede crear o desviar comercio (\autor{Jacob Viner}, 1950), y el mercado común añade las ganancias dinámicas de la competencia y la escala que se han visto en el bloque~\ref{sec:mercado-unico} \vertema{3.B.10}.
  \end{itemize}
\item \emph{Relevancia}: es la libertad más desarrollada y la que concentra la mayor parte del comercio entre los Estados (gráfico~\ref{fig:comercio-intra}). Se apoya en la unión aduanera (arts. 28-32 TFUE): sin aranceles internos y con un arancel exterior común, que gestionan las aduanas nacionales con un Código Aduanero común \vertema{3.B.41}.
\item \emph{El Derecho originario} prohíbe:
  \begin{itemize}
  \item los derechos de aduana y las \emph{exacciones de efecto equivalente} (art. 30);
  \item las restricciones cuantitativas a la importación y a la exportación y las \concepto{medidas de efecto equivalente} (arts. 34~y 35);
  \item la discriminación fiscal de los productos de otros Estados (art. 110).
  \end{itemize}
\item \emph{Excepciones} (art. 36): razones de orden público, moralidad y seguridad públicas, protección de la salud y la vida de las personas y animales, preservación de los vegetales, protección del patrimonio artístico, histórico o arqueológico y de la propiedad industrial y comercial, siempre que no sean una discriminación arbitraria ni una restricción encubierta.
\end{itemize}

\subsubsection{La jurisprudencia: de \emph{Dassonville} a \emph{Keck}}\label{sec:jurisprudencia-bienes}

\begin{itemize}
\item \emph{\concepto{Dassonville}} (asunto 8/74, 1974): es medida de efecto equivalente «toda normativa comercial de los Estados miembros que pueda obstaculizar directa o indirectamente, actual o potencialmente, el comercio intracomunitario». Una definición tan amplia que casi cualquier norma nacional podía caer en ella.
\item \emph{\concepto{Cassis de Dijon}} (asunto 120/78, 1979):\footnote{Alemania no dejaba vender como licor ninguna bebida de menos de 25~grados, con el argumento de proteger la salud y evitar que el consumidor se sintiera engañado. El licor de grosella \emph{Cassis de Dijon}, fabricado legalmente en Francia, tenía unos 15~grados. El Tribunal respondió que bastaba con informar de la graduación en la etiqueta.}
  \begin{itemize}
  \item todo producto fabricado y comercializado legalmente en un Estado debe poder venderse, en principio, en los demás: es el \concepto{principio de reconocimiento mutuo};
  \item un Estado solo puede impedirlo por las razones del art. 36~o por «\emph{exigencias imperativas}» de interés general (eficacia de los controles fiscales, protección de la salud pública, lealtad de las transacciones comerciales y defensa de los consumidores), y con medidas proporcionadas.
  \end{itemize}
\item \emph{\concepto{Keck y Mithouard}} (asuntos C-267/91 y C-268/91, 1993) puso límite a \emph{Dassonville}: las normas que regulan \emph{modalidades de venta} (horarios, venta a pérdida, publicidad) no son medidas de efecto equivalente si se aplican a todos los operadores y afectan igual, de hecho y de Derecho, a los productos nacionales y a los importados. Las que regulan el \emph{producto} (composición, envase, etiquetado) siguen sometidas a \emph{Cassis}.
\item \emph{Las fresas españolas} (Comisión contra Francia, C-265/95, 1997): Francia incumplió el Tratado al no impedir durante años que los agricultores destruyeran camiones de fruta y hortaliza españolas. Los Estados también deben actuar contra los obstáculos que crean los particulares. De ahí nació el «reglamento fresa» de 1998, que hoy forma parte del IMERA (\ref{sec:relanzamientos}).
\end{itemize}

\subsubsection{Las barreras del informe Cecchini y su desmontaje}\label{sec:barreras}

\paragraph{Barreras físicas}

\begin{itemize}
\item \emph{Qué eran}: los controles en frontera (documentos de aduana, licencias de transporte, controles sanitarios y fitosanitarios) que fragmentaban el mercado.
\item \emph{Cómo se suprimieron} (1993):
  \begin{itemize}
  \item armonizando los controles sanitarios y de seguridad y aplicando el reconocimiento mutuo de los controles hechos en origen;
  \item con un sistema nuevo de estadística del comercio (Intrastat) y de declaración del IVA, sin ajustes en frontera;
  \item reforzando la frontera exterior: las aduanas nacionales funcionan «como si fueran una sola», con el programa de cooperación Aduana 2021-2027\nocite{UE2021Aduana}, para evitar la \concepto{desviación del tráfico} (\emph{trade deflection}): que las importaciones entren por la aduana más permisiva.
  \end{itemize}
\item \emph{Lo nuevo}: el Código Aduanero reformado (Reglamento (UE) 2026/2108), que se aplica de forma gradual desde el 21~de septiembre de 2026, crea una Autoridad Aduanera de la UE con sede en Lille y un centro de datos aduaneros común; desde el 1~de julio de 2026 los pequeños envíos de menos de 150~euros pagan un derecho fijo de 3~euros por artículo y, desde el 1~de noviembre, una tasa de gestión de 2~euros\nocite{CE2026ReformaAduanera}.
\end{itemize}

\paragraph{Barreras técnicas}

\begin{itemize}
\item \emph{Qué son}: reglamentaciones distintas en cada Estado para productos parecidos. Son las más dañinas, porque obligan a las empresas a adaptar el producto a cada mercado, encarecen los precios y exigen a las administraciones verificar y certificar cada vez.
\item \emph{El enfoque tradicional} (hasta 1985): \emph{armonización detallada} de cada producto por directiva, con unanimidad. Fue lento y costoso, porque había que acordar cada detalle técnico, y los retrasos en la transposición lo empeoraban.
\item \emph{El nuevo enfoque} (Resolución del Consejo de 7~de mayo de 1985), que sigue vigente:
  \begin{itemize}
  \item las directivas fijan solo los \emph{requisitos esenciales} (de seguridad, de salud, de medio ambiente);
  \item las especificaciones técnicas las elaboran los organismos europeos de normalización (CEN, Cenelec, ETSI) como \emph{normas armonizadas} voluntarias; quien las cumple disfruta de una \emph{presunción de conformidad} y lleva el marcado CE;\footnote{Las normas armonizadas se elaboran a petición de la Comisión con la participación de la industria. No son obligatorias, pero el producto que las sigue se presume conforme con los requisitos esenciales de la directiva o del reglamento.}
  \item fuera de lo armonizado, rige el reconocimiento mutuo.
  \end{itemize}
\item \emph{Los riesgos del reconocimiento mutuo}, que explican por qué no basta con él:
  \begin{itemize}
  \item \emph{deslocalización}: las empresas podrían instalarse en el Estado con la norma más laxa y vender desde allí a todos;
  \item \emph{carrera a la baja}: los Estados podrían rebajar sus normas para atraerlas. Por eso el reconocimiento mutuo va con un mínimo armonizado.
  \end{itemize}
\item \emph{El derecho derivado vigente}\nocite{UE2019Reconocimiento,UE2015TRIS}:
  \begin{itemize}
  \item el \emph{nuevo marco legislativo} (Reglamento (CE) 765/2008 y Decisión 768/2008/CE), sobre acreditación, vigilancia del mercado y marcado CE, y el Reglamento (UE) 2019/1020 de vigilancia del mercado;
  \item el \concepto{Reglamento (UE) 2019/515 de reconocimiento mutuo}, aplicable desde el 19~de abril de 2020, que sustituyó al de 2008: el fabricante puede presentar una «declaración de reconocimiento mutuo», la autoridad que deniega el acceso debe justificarlo y los conflictos se resuelven con SOLVIT;
  \item la \emph{Directiva (UE) 2015/1535}: cada Estado notifica a la Comisión sus proyectos de reglamentos técnicos, que quedan suspendidos 3~meses para que la Comisión y los demás Estados puedan oponerse (procedimiento TRIS)\nocite{UE2015TRIS};
  \item el Reglamento (UE) 2018/302 sobre el \emph{bloqueo geográfico}, aplicable desde el 3~de diciembre de 2018, que prohíbe discriminar en las ventas en línea por la nacionalidad o la residencia del cliente;
  \item el Reglamento (UE) 2023/988 de seguridad general de los productos y el de ecodiseño (2024/1781), cuyo \emph{pasaporte digital de producto} tiene registro desde julio de 2026\nocite{CE2026Pasaporte}.
  \end{itemize}
\end{itemize}

\paragraph{Barreras fiscales: la armonización de la imposición indirecta}\label{sec:barreras-fiscales}

\begin{itemize}
\item \emph{Por qué se armoniza}: sin armonización, las diferencias entre los sistemas fiscales deciden adónde van los bienes, los servicios y los factores, y no la eficiencia. Hay 2~argumentos:
  \begin{itemize}
  \item \emph{las distorsiones}: se compra, se produce o se invierte donde el impuesto es menor, no donde el producto es mejor o el proyecto más rentable; y si el impuesto se ajusta en frontera, hacen falta controles que rompen el mercado;
  \item \emph{la competencia fiscal a la baja}: cada Estado tiene incentivos a bajar los impuestos sobre las bases móviles para atraerlas; si todos lo hacen, ninguno las atrae y todos recaudan menos. Es un dilema del prisionero, como la carrera de subvenciones (\ref{sec:ayudas}), que la coordinación resuelve \vertema{3.A.14}.
  \end{itemize}
\item \emph{Por qué avanza poco}: la fiscalidad no es una competencia exclusiva de los Estados en el sentido del TFUE, pero los Estados conservan su soberanía fiscal y la armonización exige \emph{unanimidad} en el Consejo (arts. 113~y 115~TFUE). Los Estados con impuestos bajos temen perder su ventaja, y los de impuestos altos, su recaudación. Por eso la Unión no pretende uniformar los sistemas, sino hacerlos compatibles con el mercado interior, con directivas y con los principios de subsidiariedad y proporcionalidad.
\item \emph{Las 3~estrategias posibles}, de más a menos ambiciosa:
  \begin{itemize}
  \item la \emph{igualación}: un mismo sistema en todos los Estados; elimina las distorsiones, pero es desproporcionada y políticamente inviable;
  \item la \emph{aproximación}: bases comunes y horquillas de tipos solo en los impuestos que estorban al mercado; es la que ha seguido la Unión, de forma \emph{explícita} (directivas) en el IVA y los impuestos especiales y \emph{implícita} (códigos de conducta, recomendaciones) en los impuestos directos;
  \item la \emph{diferencial}: armonizar lo que maximice el bienestar del conjunto; atractiva en teoría, pero imposible de medir en la práctica.
  \end{itemize}
\item \emph{2 principios para gravar el comercio entre Estados}:\footnote{Los 2~principios, en detalle:
\begin{itemize}
\item con \emph{origen}, cada bien soporta el impuesto del país donde se produce, y la recaudación es de ese país; para que 2~bienes de 2~países compitan en igualdad dentro de un mismo mercado, hay que armonizar bases y tipos;
\item con \emph{destino}, el bien soporta el impuesto del país donde se consume, y la recaudación es del importador; cada Estado conserva su tipo, pero hay que controlar las operaciones.
\end{itemize}
Los Estados se han resistido al origen porque les quita autonomía y porque los importadores netos perderían recaudación.}
  \begin{itemize}
  \item \emph{origen}: no exige controles, pero sí armonizar los tipos;
  \item \emph{destino}: deja a cada Estado su tipo, pero exige controlar las operaciones entre Estados.
  \end{itemize}
\item \emph{La armonización del IVA}, en 3~etapas, de las que se han cumplido 2\nocite{UE2006IVA,UE2022Tipos}:
  \begin{itemize}
  \item \emph{implantar el IVA en todos los Estados}: las 2~primeras directivas de 1967 sustituyeron los impuestos en cascada por el IVA, que desde entonces es requisito para la adhesión (España lo implantó en 1986) \vertema{4.B.20};
  \item \emph{unificar la base imponible}: la Sexta Directiva (1977), refundida en la Directiva 2006/112/CE, hizo que el impuesto gravara las mismas operaciones en todos los Estados; una parte de esa base es, desde 1970, un recurso propio de la Unión \vertema{3.B.37};
  \item \emph{aproximar los tipos}: es la etapa que falta; el \emph{paquete Cockfield} (1987) propuso pasar al origen con tipos aproximados y no prosperó.
  \end{itemize}
\item \emph{El régimen «transitorio» de 1993}, que sigue vigente:
  \begin{itemize}
  \item las ventas entre empresas de distintos Estados están exentas en origen y tributan en destino, sin ajustes en frontera;
  \item el control se hace con la contabilidad de las empresas y el intercambio de información entre administraciones (sistema VIES);
  \item el «régimen definitivo» en origen nunca llegó, y la Comisión se ha centrado en mejorar el transitorio: los servicios entre empresas tributan en destino por \emph{inversión del sujeto pasivo} desde 2010, y las ventas a distancia y digitales a consumidores, con la ventanilla única (OSS) desde 2021\nocite{UE2006IVA}.
  \end{itemize}
\item \emph{Los tipos}: la Directiva (UE) 2022/542 permite un tipo general no inferior al 15~\%, hasta 2~tipos reducidos no inferiores al 5~\% y, para una lista limitada de bienes, un tipo inferior al 5~\% y la exención con derecho a deducción.\footnote{En España, el tipo general es el 21~\%, el reducido el 10~\% y el superreducido el 4~\% \vertema{4.B.20}.}
\item \emph{El IVA en la era digital} (paquete ViDA, Directiva (UE) 2025/516)\nocite{CE2026ViDA}:
  \begin{itemize}
  \item registro único en la Unión y plataformas de alojamiento de corta duración y de transporte de viajeros como sujetos pasivos, desde julio de 2028\nocite{CE2026ViDA};
  \item facturación electrónica y declaración digital de las operaciones entre empresas de distintos Estados, desde julio de 2030, para atajar el fraude\nocite{CE2026ViDA}.
  \end{itemize}
\item \emph{Los impuestos especiales}, con el mismo esquema que el IVA:
  \begin{itemize}
  \item la Directiva de 1992 armonizó la estructura del alcohol, el tabaco y los hidrocarburos, con tipos mínimos y el principio de destino;
  \item la Directiva 2003/96/CE extendió los mínimos a la electricidad, el gas natural y el carbón;
  \item la Directiva (UE) 2020/262, aplicable desde febrero de 2023, renovó el régimen general, y la 2020/1151 puso al día la del alcohol\nocite{UE2020ImpuestosEspeciales};
  \item las reformas de la fiscalidad de la energía (propuesta de 2021) y del tabaco (2025) siguen sin la unanimidad necesaria\nocite{CE2026Energia}.
  \end{itemize}
\item \emph{El instrumento de cooperación}: el programa \emph{Fiscalis} (2021-2027) financia la cooperación entre las administraciones tributarias y sus sistemas de información\nocite{UE2021Fiscalis}.
\item \emph{Balance}: no puede hablarse de un mercado único en términos fiscales. El IVA transitorio es complejo y vulnerable al fraude, y las diferencias de tipos de los impuestos especiales alientan el contrabando. La imposición directa, que afecta más a los capitales, se trata en \ref{sec:capitales-marco}.
\end{itemize}

\begin{anotaciones}
\textbf{La imposición indirecta en la Unión, en cifras.}
\begin{itemize}
\item \emph{Tipos generales del IVA} (2026): del 17~\% de Luxemburgo al 27~\% de Hungría; España, el 21~\%. El mínimo es el 15~\%\nocite{CE2026TiposIVA}.
\item \emph{La brecha del IVA}: en 2023 los Estados dejaron de recaudar 128.000~millones de euros, el 9,5~\% de lo que debían ingresar, por fraude, evasión, quiebras y errores; la brecha había bajado hasta 2022 y volvió a subir\nocite{CE2025BrechaIVA}.
\item \emph{El fraude carrusel}: la exención de las ventas entre Estados permite a empresas pantalla comprar sin IVA, venderlo cobrado y desaparecer sin ingresarlo; la Comisión lo estimaba en unos 50.000~millones de euros al año cuando propuso el régimen definitivo\nocite{CE2017RegimenDefinitivo}.
\end{itemize}
\end{anotaciones}

\paragraph{Restricciones de acceso a los contratos públicos}

\begin{itemize}
\item \emph{Qué son}: procedimientos de compra pública que favorecen a las empresas nacionales y segmentan el mercado. Más de 250.000~poderes adjudicadores gastan cada año unos 2,6~billones de euros, el 15~\% del PIB de la Unión\nocite{CE2026Contratacion}.
\item \emph{Por qué abrirlos}: tiene un \emph{efecto directo}, bajar los precios que paga el sector público con más licitadores, y un \emph{efecto de reestructuración}, más dinámico: concentraciones que racionalizan la producción y reducen los costes de desarrollo, como en el modelo de \ref{sec:modelo-integracion}.
\item \emph{Cómo se han abierto}:
  \begin{itemize}
  \item las directivas de 2014 sobre contratación pública y concesiones (2014/23/UE, 2014/24/UE y 2014/25/UE, aplicables desde 2016)\nocite{CE2026Contratacion}, que simplifican los procedimientos y extienden las reglas a los sectores del agua, la energía, el transporte y los servicios postales;
  \item la publicación de las licitaciones en el \emph{Diario Oficial} y los recursos frente a la discriminación.
  \end{itemize}
\item \emph{Lo pendiente}: la adjudicación directa a empresas de otros Estados sigue siendo muy poco frecuente, y la estrategia de 2025 anuncia una revisión de las directivas para simplificarlas.
\end{itemize}

\subsubsection{Valoración}\label{sec:bienes-valoracion}

\begin{itemize}
\item[(+)] El comercio de bienes dentro de la Unión creció más que el PIB hasta 2008 sin desviar apenas comercio del resto del mundo; los precios convergieron y los márgenes de la industria bajaron.
\item[(+)] Para los Estados del Este, el comercio con la Unión ha sido el gran motor de su integración: sus importaciones de otros Estados de la Unión rondan el 50~\% de su PIB.
\item[(--)] En la Europa de los 15~la integración se ha estancado desde mediados de los 2000, y la principal fuerza del comercio intracomunitario ha sido la ampliación al Este, no la profundización.
\item[(--)] Siguen las preferencias por lo nacional (\emph{home bias}) y las barreras técnicas: el 60~\% de las barreras actuales son las mismas que hace 20~años.
\end{itemize}

% --------------------------------------------------------------------------
\subsection{Mercado interior de servicios}\label{sec:servicios}

\subsubsection{Justificación y relevancia}\label{sec:servicios-relevancia}

\begin{itemize}
\item La lógica económica es la misma que en los bienes, pero el peso es mayor: los servicios suponen en torno al 75~\% del PIB de la Unión\nocite{CE2025EstrategiaFicha} y solo el 8~\% del PIB se comercia en servicios dentro de ella (gráfico~\ref{fig:comercio-intra}).
\item \emph{Su peso real en el comercio es mayor del que se ve}: muchos servicios van incorporados en los bienes (diseño, I+D, ingeniería, logística, posventa). Medido en valor añadido, el comercio de servicios ronda la mitad del total, frente a menos de la cuarta parte en términos brutos \vertema{3.B.34}.\footnote{Es la \emph{servicificación} de la manufactura, que recoge la base de datos de comercio en valor añadido (TiVA) de la OCDE y la OMC. La \emph{curva de la sonrisa} de \autor{Stan Shih} (1992) la ilustra: las fases que más valor añaden son las del principio (I+D, diseño) y las del final (distribución, marca, posventa); la fabricación, en el medio, añade menos.}
\item \emph{Son más difíciles de integrar}: su comercio exige a menudo que se mueva el proveedor o el consumidor, y están muy regulados (profesiones colegiadas, licencias, requisitos de establecimiento).
\end{itemize}

\subsubsection{Marco jurídico}\label{sec:servicios-marco}

\begin{itemize}
\item \emph{2 libertades del Tratado}:\nocite{PE2026FichaServicios}
  \begin{itemize}
  \item la \concepto{libertad de establecimiento} (art. 49~TFUE): un nacional o una empresa de un Estado puede instalarse de forma estable en otro (una filial, una sucursal); equivale al modo 3~de prestación del Acuerdo General sobre el Comercio de Servicios \vertema{3.B.34};
  \item la \concepto{libre prestación de servicios} (art. 56~TFUE): el proveedor establecido en un Estado presta el servicio en otro sin instalarse en él (los modos 1, 2~y 4).
  \end{itemize}
\item \emph{Excepciones}: las actividades que impliquen el ejercicio del poder público (art. 51) y las medidas justificadas por razones de orden público, seguridad pública o salud pública (art. 52, aplicable a los servicios por el art. 62). Los transportes se rigen por su propio título y los servicios financieros ligados a movimientos de capitales, en armonía con la libertad de capitales.
\item \emph{La jurisprudencia}: desde \emph{Van Binsbergen} (1974), los requisitos nacionales a quien presta servicios desde otro Estado deben estar justificados por razones de interés general y ser proporcionados. \emph{Laval} y \emph{Viking} (2007) limitaron el uso de huelgas para imponer convenios a empresas de otros Estados.
\end{itemize}

\subsubsection{Implementación}\label{sec:servicios-implementacion}

\begin{itemize}
\item \emph{El Libro Blanco de 1985} apenas concretó medidas para los servicios: si para los bienes fue la culminación, para los servicios fue solo el comienzo.
\item \emph{Los años 90, enfoque sectorial}: se liberalizaron, de forma gradual, los sectores en monopolio público (telecomunicaciones, transporte aéreo, electricidad y gas, servicios postales, ferrocarril), con efectos procompetitivos notables en precios (\ref{sec:promocion}).
\item \emph{2006, enfoque horizontal: la \concepto{Directiva de servicios}} (2006/123/CE, la «Directiva Bolkestein»), transpuesta en España por las Leyes 17/2009 («ley paraguas») y 25/2009 («ley ómnibus»)\nocite{UE2006Servicios}:
  \begin{itemize}
  \item se aplica a una gran parte de los servicios (distribución, turismo, construcción, servicios a empresas, profesiones), pero excluye los financieros, las comunicaciones electrónicas, el transporte, la sanidad, el trabajo temporal, los servicios sociales y los audiovisuales;
  \item obliga a revisar los regímenes de autorización, que deben ser no discriminatorios, necesarios y proporcionados, y prohíbe algunos requisitos (de nacionalidad o pruebas de necesidad económica);
  \item crea \emph{ventanillas únicas} y la cooperación administrativa entre Estados (Sistema de Información del Mercado Interior);
  \item el principio del \emph{país de origen} de la propuesta de 2004, que hacía cumplir solo las normas del Estado de establecimiento, se sustituyó por una «libertad de prestación de servicios» con excepciones;\footnote{La propuesta inicial de \autor{Frits Bolkestein} generó una fuerte oposición social (el «fontanero polaco») y se rebajó en el Parlamento Europeo. Tampoco salió adelante la reforma de los colegios profesionales.}
  \item su valoración es positiva, pero, al ser una directiva, su efecto depende de cómo la transpone y aplica cada Estado.
  \end{itemize}
\item \emph{Las cualificaciones profesionales}: la Directiva 2005/36/CE, reformada por la 2013/55/UE, regula el reconocimiento de títulos:
  \begin{itemize}
  \item \emph{automático} para 7~profesiones sectoriales: médicos, enfermeros, odontólogos, veterinarios, matronas, farmacéuticos y arquitectos;
  \item \emph{general} para las demás, con posibles medidas compensatorias (periodo de prácticas o prueba de aptitud);
  \item la \emph{tarjeta profesional europea} electrónica, desde 2016, para enfermeros, farmacéuticos, fisioterapeutas, guías de montaña y agentes inmobiliarios;
  \item la Directiva (UE) 2018/958 obliga a hacer un \emph{test de proporcionalidad} antes de regular una profesión.\footnote{Las profesiones reguladas emplean en torno al 22~\% de los trabajadores de la Unión. El paquete de servicios de 2017 propuso además una tarjeta electrónica de servicios para facilitar los trámites de las empresas de construcción y de servicios a empresas, que el Parlamento Europeo rechazó.}
  \end{itemize}
  En septiembre de 2026 la Comisión propuso reformarla para hacer el reconocimiento digital y más rápido\nocite{CE2026Movilidad}.
\item \emph{Los servicios financieros}: el \emph{pasaporte único} permite a un banco, una aseguradora o una empresa de inversión autorizada en un Estado operar en todos, con el código normativo único (requisitos de capital, recuperación y resolución, garantía de depósitos), MiFID II y Solvencia II, y la zona única de pagos en euros (SEPA) \vertema{3.B.44}.
\item \emph{El mercado único digital} (estrategia de 2015): portabilidad de los contenidos en línea, fin del \emph{roaming} y, desde 2022, el \concepto{Reglamento de Servicios Digitales} (DSA, 2022/2065) y el \concepto{Reglamento de Mercados Digitales} (DMA, 2022/1925)\nocite{UE2022DSA,UE2022DMA}:
  \begin{itemize}
  \item la DSA ordena la responsabilidad de los intermediarios en línea por los contenidos ilícitos, con obligaciones mayores para las 25~plataformas y 3~buscadores muy grandes (más de 45~millones de usuarios); sus primeras multas son a X (120~millones de euros, diciembre de 2025), Temu (200~millones) y AliExpress (550~millones, 2026)\nocite{CE2026DSA};
  \item la DMA regula \emph{ex ante} a los guardianes de acceso y se cuenta con la competencia (\ref{sec:abuso}).
  \end{itemize}
\end{itemize}

\subsubsection{Valoración y retos}\label{sec:servicios-valoracion}

\begin{itemize}
\item[(--)] Los servicios son la gran asignatura pendiente: un arancel equivalente de hasta el 110~\% según el FMI\nocite{FMI2024REO}. Las barreras más altas están en la construcción, los servicios profesionales y la distribución, según los propios estudios de la Comisión.
\item[(--)] La fragmentación frena la productividad de los servicios europeos, que crece menos que la de Estados Unidos. Las telecomunicaciones son el ejemplo: hay decenas de operadores en la Unión, pero en cada país el consumidor elige entre 3~o 4; tienen los inconvenientes del oligopolio sin las economías de escala de sus rivales estadounidenses o chinos.
\item \emph{Retos}: la regulación de las plataformas y de la economía colaborativa, que cada Estado regula a su manera;\footnote{Plataformas digitales que ponen en contacto a particulares para alquilar, compartir o prestar servicios (alojamiento, transporte, trabajos puntuales). Ofrecen precios más bajos y más variedad, pero plantean problemas de protección del consumidor, derechos de los trabajadores y fiscalidad que cada Estado resuelve de forma distinta. La Directiva (UE) 2024/2831 regula el trabajo en plataformas digitales.} la armonización de las autorizaciones de servicios y las restricciones territoriales de suministro, 2~de los «\emph{Terrible Ten}».
\end{itemize}

% --------------------------------------------------------------------------
\subsection{Libre circulación de personas}\label{sec:personas}

\subsubsection{Justificación}\label{sec:personas-justificacion}

\begin{itemize}
\item La constitución del mercado único no solo pretende la libre circulación de bienes y servicios, sino también la de los factores de producción.
\item \emph{Teoría neoclásica}: si el trabajo y el capital se mueven hacia donde son más productivos, se agotan las ganancias del intercambio y la relación marginal de sustitución técnica se iguala a la relación de precios de los factores en todas partes, $\mathrm{RMST}=w/r$ \vertema{3.B.18}. En la práctica, cuando cada trabajador encuentra el empleo que mejor se ajusta a sus capacidades, es más productivo y gana más.
\item \emph{Geografía económica}: la movilidad del trabajo cualificado es una fuerza de aglomeración que puede vaciar la periferia \vertema{3.B.10}.
\item \emph{Evidencia de la inmigración}:
  \begin{itemize}
  \item en el país de destino, el efecto sobre los salarios y el paro de los nativos es pequeño, como mostró \textcite{Card1990Impact} con el éxodo de Mariel;\footnote{\autor{David Card} recibió el Premio Nobel de Economía en 2021 «por sus contribuciones empíricas a la economía laboral»\nocite{Nobel2021CardAngristImbens}. En 1980, unos 125.000~cubanos salieron del puerto de Mariel hacia Estados Unidos, y la población activa de Miami creció en torno a un 7~\%. \autor{Card} comparó los salarios y el paro de Miami con los de ciudades parecidas (Atlanta, Houston, Los Ángeles, Tampa) y apenas encontró diferencias. Una explicación es que los nativos y el capital también se mueven entre ciudades, y diluyen el efecto.}
  \item en el país de origen, ganancias por las remesas y pérdidas por la fuga de cerebros;
  \item para los propios migrantes, ganancias claras.
  \end{itemize}
\end{itemize}

\subsubsection{Marco jurídico e implementación}\label{sec:personas-marco}

\begin{itemize}
\item \emph{El Tratado de Roma} (1957) ya reconocía la libre circulación de \emph{trabajadores}. Hoy, el art. 45~TFUE suprime toda discriminación por nacionalidad en el empleo, la retribución y las condiciones de trabajo, y da derecho a:
  \begin{itemize}
  \item responder a ofertas efectivas de trabajo;
  \item desplazarse libremente con ese fin por el territorio de los Estados;
  \item residir en un Estado para trabajar en él, y quedarse después.
  \end{itemize}
  Admite límites por razones de orden público, seguridad y salud públicas, y excluye los empleos en la administración pública que impliquen ejercicio de autoridad (como el de técnico comercial y economista del Estado). El Reglamento (UE) 492/2011 la desarrolla.
\item \emph{Maastricht} (1992) creó la \concepto{ciudadanía de la Unión}, que da a todo nacional de un Estado el derecho a circular y residir libremente, sea o no trabajador (art. 21~TFUE).\footnote{Los tratados de adhesión de 2003 y 2005 permitieron a los Estados anteriores aplicar hasta 7~años de periodos transitorios a la libre circulación de los trabajadores de los nuevos Estados. España los aplicó a los de Rumanía.}
\item \emph{La Directiva 2004/38/CE} regula la residencia:
  \begin{itemize}
  \item \emph{hasta 3~meses}: basta un documento de identidad o un pasaporte válidos;
  \item \emph{más de 3~meses}: hay que trabajar o tener recursos suficientes y un seguro de enfermedad, para no ser una carga para la asistencia social del Estado de acogida; se puede exigir el registro;
  \item \emph{residencia permanente}: tras 5~años de residencia legal ininterrumpida.
  \end{itemize}
\item \emph{El espacio Schengen}: el Acuerdo de Schengen (1985) y su Convenio (1990), intergubernamentales, se aplicaron desde 1995 y el Tratado de Ámsterdam (1997) los integró en el Derecho de la Unión\nocite{CE2026Schengen}:
  \begin{itemize}
  \item suprime los controles en las fronteras interiores para todas las personas, aunque pueden restablecerse de forma temporal ante amenazas graves (el Código de fronteras Schengen reformado en 2024 los limita a 2~años, prorrogables 1~más); en noviembre de 2026 los mantenían 7~Estados de la Unión (Alemania, Austria, Dinamarca, Francia, Italia, Polonia y Suecia) y Noruega;
  \item armoniza los controles en la frontera exterior y la política de visados de corta duración, ahora con el sistema de entradas y salidas, plenamente operativo desde abril de 2026;
  \item crea la cooperación policial y judicial (el Sistema de Información de Schengen);
  \item tiene 29~miembros: 25~Estados de la Unión (Bulgaria y Rumanía, de pleno derecho desde el 1~de enero de 2025) más Islandia, Liechtenstein, Noruega y Suiza. Fuera quedan Irlanda, que tiene una exención, y Chipre, pendiente de la decisión del Consejo.\footnote{Mónaco, San Marino y el Vaticano no son miembros, pero tienen fronteras abiertas con el espacio Schengen. El acuerdo entre la Unión y el Reino Unido sobre Gibraltar, firmado el 14~de julio de 2026 y aplicado de forma provisional desde el día siguiente, suprime los controles de personas y mercancías entre Gibraltar y España, con controles de Schengen en su puerto y su aeropuerto; falta la aprobación del Parlamento Europeo.}
  \end{itemize}
\end{itemize}

\begin{anotaciones}
\textbf{El mapa del espacio Schengen.} El gráfico~\ref{fig:mapa-schengen} distingue los 25~Estados de la Unión que aplican el acervo de Schengen, los 4~asociados de fuera de la Unión y los 2~Estados de la Unión que no lo aplican del todo, y marca con trama los que mantenían controles interiores temporales en noviembre de 2026.
\begin{figuracaja}
\centering
\caption{El espacio Schengen, noviembre de 2026}\label{fig:mapa-schengen}
\grafico[0.92\textwidth]{mapa-schengen}
\fuente{elaboración propia a partir de la Comisión Europea (miembros y controles interiores notificados, consulta del 19/11/2026)\nocite{CE2026Schengen,CE2026SchengenMiembros}; contornos de los países de Natural Earth\nocite{NaturalEarth2026}. No se dibujan las regiones ultraperiféricas (las Canarias, Madeira y las Azores sí están en Schengen). Chipre participa en la cooperación de Schengen, pero el Consejo aún no ha suprimido sus controles interiores.}
\end{figuracaja}
\end{anotaciones}

\begin{itemize}
\item \emph{El desplazamiento de trabajadores}: un \concepto{trabajador desplazado} es el que su empresa envía temporalmente a otro Estado a prestar un servicio, y sigue cotizando en el de origen\nocite{UE2018Desplazamiento}.
  \begin{itemize}
  \item La Directiva 96/71/CE le garantizaba las condiciones mínimas del Estado de acogida; la Directiva (UE) 2018/957, aplicable desde julio de 2020, cambió el «salario mínimo» por la \emph{misma remuneración} que los locales y, si el desplazamiento supera 12~meses (18~con notificación motivada), le aplica casi todo el Derecho laboral del Estado de acogida.
  \item El objetivo es evitar el \concepto{dumping social}: que las empresas ganen cuota en otros Estados solo por pagar menos o cotizar menos.
  \item Su volumen: 3,6~millones de documentos A1 de desplazamiento en 2024\nocite{CE2025A1}. Una declaración electrónica común, primer resultado de la hoja de ruta de 2026, simplificará el trámite.
  \end{itemize}
\item \emph{Instrumentos} que facilitan la movilidad:
  \begin{itemize}
  \item la red EURES, que pone en común las ofertas de empleo de los servicios públicos de empleo;
  \item el programa Erasmus+, de movilidad de estudiantes, profesores y trabajadores en formación;
  \item la \emph{Autoridad Laboral Europea} (Bratislava, Reglamento (UE) 2019/1149), que coordina las inspecciones y la aplicación de las normas de movilidad laboral desde octubre de 2019\nocite{UE2019ALE}.
  \end{itemize}
\end{itemize}

\subsubsection{Valoración y retos}\label{sec:personas-valoracion}

\begin{itemize}
\item \emph{Indicadores}: unos 14~millones de ciudadanos de la Unión viven o trabajan de forma duradera en otro Estado, a los que se suman 1,9~millones de trabajadores fronterizos y 3,6~millones de desplazados\nocite{CE2026Movilidad}. Los movimientos van sobre todo del Este y el Sur hacia los Estados de salarios más altos, y la cualificación de los trabajadores influye en su destino.
\item[(--)] \emph{La movilidad es baja}: mucho menor que entre los estados de Estados Unidos, lo que la hace un mal mecanismo de ajuste ante perturbaciones asimétricas en la zona del euro \vertema{3.B.16}.
\item \emph{Barreras} que frenan la movilidad:
  \begin{itemize}
  \item la lengua y la cultura;
  \item el reconocimiento de títulos;
  \item la falta de coordinación de la fiscalidad del trabajo y de la seguridad social (pensiones, desempleo).\footnote{Los Estados siguen siendo responsables de sus sistemas educativos. El suplemento europeo al título y los créditos ECTS del proceso de Bolonia ayudan al reconocimiento, pero no lo garantizan.}
  \end{itemize}
\item \emph{Retos}: el reconocimiento de cualificaciones, también las de nacionales de terceros países, y la portabilidad de los derechos sociales, que el paquete de movilidad laboral de septiembre de 2026 quiere digitalizar\nocite{CE2026Movilidad}.
\end{itemize}

% --------------------------------------------------------------------------
\subsection{Libre circulación de capitales}\label{sec:capitales}

\subsubsection{Justificación y relevancia}\label{sec:capitales-justificacion}

\begin{itemize}
\item \emph{La teoría}: la libre circulación de capitales permite:
  \begin{itemize}
  \item suavizar el consumo y financiar la inversión con ahorro ajeno, según el enfoque intertemporal de la balanza de pagos \vertema{3.B.12};
  \item diversificar el riesgo, según la teoría de la cartera \vertema{3.B.18};
  \item transferir tecnología y gestión con la inversión directa \vertema{3.B.18}.
  \end{itemize}
\item \emph{En una unión monetaria es imprescindible}: los ataques al Sistema Monetario Europeo de 1992-1993 mostraron que, con capitales libres, los tipos de cambio fijos pero ajustables eran frágiles \vertema{3.B.43}.
\item \emph{Relevancia}: es la más amplia de las 4~libertades, porque es la única que se extiende a terceros países (\emph{erga omnes}).
\end{itemize}

\subsubsection{Implementación y marco jurídico}\label{sec:capitales-marco}

\begin{itemize}
\item \emph{De Roma al mercado único}:
  \begin{itemize}
  \item el Tratado de Roma solo pedía liberalizar los capitales «en la medida necesaria» para el mercado común;
  \item las directivas de 1960 y 1962 liberalizaron la inversión directa, los créditos comerciales y las acciones cotizadas, y las de 1985-1986, los títulos a largo plazo y los no cotizados;
  \item la \concepto{Directiva 88/361/CEE} fijó la \emph{liberalización completa} el 1~de julio de 1990, con una cláusula de salvaguardia ante movimientos a corto plazo excepcionales y transición hasta finales de 1992 para España, Irlanda, Grecia y Portugal, prorrogable para los 2~últimos\nocite{UE1988Capitales};
  \item el \emph{Tratado de Maastricht} la llevó al Tratado, desde el 1~de enero de 1994.
  \end{itemize}
\item \emph{El TFUE}:
  \begin{itemize}
  \item el art. 63~prohíbe toda restricción a los movimientos de capitales y a los pagos entre Estados y con terceros países;
  \item el art. 64~mantiene las restricciones frente a terceros países vigentes el 31~de diciembre de 1993, y el 66~permite salvaguardias frente a ellos de hasta 6~meses;
  \item el art. 65~admite medidas fiscales que distingan por residencia, medidas de supervisión prudencial, declaraciones a efectos estadísticos y medidas de orden o seguridad públicos. Por esa vía, Chipre (2013) y Grecia (2015) impusieron controles de capitales durante sus crisis, levantados por completo en 2015 y 2019 \vertema{3.B.15};\footnote{Además, los arts. 143~y 144~TFUE permiten medidas de salvaguardia de la balanza de pagos a los Estados que no tienen el euro, y los arts. 75~y 215, sanciones financieras contra el terrorismo y en el marco de la política exterior.}
  \item las \emph{acciones de oro}: el Tribunal de Justicia condenó a España (C-463/00, 2003) por el régimen de autorización previa de las empresas privatizadas (Repsol, Telefónica, Endesa…), que restringía la libre circulación de capitales.
  \end{itemize}
\item \emph{El control de las inversiones extranjeras directas}: la libertad frente a terceros países tiene un límite de seguridad. El Reglamento (UE) 2019/452 creó un marco de cooperación entre los Estados que las examinan; el nuevo Reglamento (UE) 2026/1386, en vigor desde julio de 2026 y aplicable desde enero de 2028, obliga a todos los Estados a tener un mecanismo de control y fija unos sectores mínimos (semiconductores, inteligencia artificial, materias primas críticas)\nocite{UE2026Inversiones} \vertema{3.B.41}.
\item \emph{La Unión de Mercados de Capitales}: lanzada en 2015 para que la financiación de las empresas dependa menos de los bancos, con pocos avances, se relanzó en marzo de 2025 como \concepto{Unión de Ahorros e Inversiones}: cuentas de ahorro e inversión, titulización (acuerdo político en octubre de 2026) y la supervisión directa por la Autoridad Europea de Valores y Mercados de las grandes infraestructuras, en negociación\nocite{CE2025UAI}. Aquí no se trata su diseño \vertema{3.B.44}.
\item \emph{La armonización de la imposición directa}, que exige unanimidad (art. 115~TFUE). Afecta más a los capitales y a la libertad de establecimiento que a los bienes: las diferencias en el impuesto sobre sociedades pueden llevar los beneficios y las sedes a donde menos se paga, no a donde más rinden, y alimentar una competencia fiscal a la baja \vertema{4.B.18}. Por eso llegó más tarde y se ha limitado a corregir distorsiones concretas, en 4~frentes:
  \begin{itemize}
  \item \emph{evitar la doble imposición dentro de los grupos}: la Directiva de matrices y filiales (1990, hoy 2011/96/UE) exime de gravar los dividendos que una filial paga a su matriz de otro Estado; la de intereses y cánones (2003/49/CE) hace lo mismo con esos pagos; y la de fusiones (1990, hoy 2009/133/CE) aplaza la tributación de las plusvalías de las reestructuraciones entre Estados;
  \item \emph{frenar la competencia perjudicial}: el código de conducta sobre la fiscalidad de las empresas (1997), no vinculante, que ha desmontado centenares de regímenes preferentes, y la lista de la Unión de jurisdicciones no cooperadoras (desde 2017);
  \item \emph{combatir la elusión}: la Directiva contra la elusión fiscal (ATAD, 2016/1164, ampliada en 2017), que aplica en la Unión los resultados del proyecto BEPS de la OCDE (límite a la deducción de intereses, imposición de salida, transparencia fiscal internacional y cláusula general antiabuso);
  \item \emph{fijar un mínimo}: la \emph{Directiva (UE) 2022/2523}, que aplica desde 2024 el impuesto mínimo global del 15~\% del pilar 2~de la OCDE a los grupos que facturan al menos 750~millones de euros; en enero de 2026 el Marco Inclusivo de la OCDE acordó un régimen «paralelo» (\emph{side-by-side}) que exime de parte de sus reglas a los grupos con matriz en Estados Unidos\nocite{UE2022Pilar2,OCDE2026SideBySide}.
  \end{itemize}
\item \emph{Lo que sigue pendiente} en la imposición directa:
  \begin{itemize}
  \item una base imponible común del impuesto sobre sociedades: la propuesta de base consolidada (BICCIS, 2011 y 2016) se retiró y la que la sustituye, BEFIT (2023), sigue sin acuerdo;
  \item los tipos, que van del 9~\% de Hungría a más del 30~\% de Alemania (España, el 25~\%)\nocite{TaxFoundation2026Sociedades}, y la propuesta contra las sociedades pantalla, retirada en 2025.
  \end{itemize}
\item \emph{La fiscalidad del ahorro}: con la libertad de capitales, el ahorro puede irse a donde sus rendimientos tributan menos o no se declaran. La respuesta no ha sido armonizar los tipos, sino la información:
  \begin{itemize}
  \item la Directiva de fiscalidad del ahorro (2003/48/CE) obligó a intercambiar información sobre los intereses pagados a no residentes, con una retención transitoria en Austria, Bélgica y Luxemburgo, que mantenían el secreto bancario;
  \item la \emph{directiva de cooperación administrativa} (DAC, 2011/16/UE) la sustituyó en 2016 con el intercambio automático de las cuentas financieras según el estándar común de la OCDE, y sus reformas lo han extendido a los acuerdos fiscales, la planificación agresiva, las plataformas digitales, los criptoactivos (desde 2026) y el impuesto mínimo;
  \item la Directiva (UE) 2025/50 (FASTER) agilizará, desde 2030, la devolución de las retenciones en origen a los inversores de otros Estados.
  \end{itemize}
\end{itemize}

\begin{anotaciones}
\textbf{La paradoja de \autor{Feldstein} y \autor{Horioka}.} \textcite{FeldsteinHorioka1980Domestic} observaron que, en los países de la OCDE, el ahorro y la inversión nacionales estaban muy correlacionados, lo que llevaba a saldos por cuenta corriente pequeños: como si el capital no se moviera entre países. Contradice la teoría, según la cual el ahorro debería ir a los mejores proyectos del mundo. Las explicaciones son la información asimétrica, el riesgo de tipo de cambio y regulatorio, los costes de transacción y la preferencia por lo nacional (\emph{home bias}). Con la liberalización, la correlación cayó, sobre todo dentro de la zona del euro\nocite{BlanchardGiavazzi2002Current}, y aparecieron los grandes déficits del Sur antes de 2008 \vertema{3.B.12}.
\end{anotaciones}

\subsubsection{Valoración y retos}\label{sec:capitales-valoracion}

\begin{itemize}
\item[(+)] La inversión directa entre Estados ha crecido más que la procedente de fuera de la Unión: los fundadores (Alemania, Francia, Países Bajos) son inversores netos y los Estados del Este, receptores.
\item[(--)] La inversión en cartera sigue sesgada hacia lo nacional, y la integración financiera retrocedió tras la crisis de 2008: la fragmentación hizo que empresas iguales pagaran tipos distintos según su país \vertema{3.B.44}.
\item[(--)] Una integración mal orientada puede ser costosa: \textcite{GopinathEtAl2017Capital} muestran que, tras el euro, el capital barato del Sur de Europa fue a empresas con más patrimonio pero no más productivas, y la productividad cayó.
\item \emph{Retos}: la Unión de Ahorros e Inversiones, la armonización fiscal y el equilibrio entre apertura y seguridad económica.
\end{itemize}

\begin{itemize}
\item \emph{Recapitulación del bloque~\ref{sec:libertades}}: las 4~libertades tienen la misma arquitectura (Tratado, jurisprudencia, derecho derivado) y distinto grado de avance: mucho en bienes y capitales, poco en servicios y personas (figura~\ref{figu:sintesis-libertades}).
\item \emph{Los retos comunes del mercado único}, que recoge la estrategia de 2025:
  \begin{itemize}
  \item aplicar lo que ya existe: buena parte de las barreras son normas nacionales que contradicen las europeas o que se transponen tarde;
  \item completar los servicios, con menos profesiones reguladas y autorizaciones más sencillas;
  \item integrar los mercados de capitales y la fiscalidad de las empresas, con el 28.º régimen y la Unión de Ahorros e Inversiones;
  \item preservar el mercado en las crisis sin cierres de fronteras, con el IMERA y el control de las inversiones extranjeras.
  \end{itemize}
\item Pero un mercado sin barreras no es aún un mercado competitivo: hace falta que las empresas y los Estados no levanten otras. Es la tarea de la política de competencia.
\end{itemize}

\begin{infografia}[H]
\centering
\caption{Las 4~libertades de un vistazo: mismo patrón, distinto avance}\label{figu:sintesis-libertades}
\grafico[\textwidth]{sintesis-libertades}
\fuente{elaboración propia. El grado de avance es una valoración cualitativa a partir del bloque~\ref{sec:libertades} y de las estimaciones del FMI\nocite{FMI2024REO} y de la Comisión\nocite{CE2025EstrategiaFicha}.}
\end{infografia}

% ==========================================================================
\section{La política europea de competencia}\label{sec:competencia}
\minutos{10}

\begin{itemize}
\item Una vez vistas las libertades, pasamos a las reglas que impiden que las empresas o los Estados vuelvan a cerrar el mercado.
\item \emph{Conclusión anticipada}: la política de competencia es la otra cara del mercado único. Controla a las empresas (acuerdos, abusos, concentraciones) y, cosa única en el mundo, a los Estados (ayudas). Su criterio es el bienestar del consumidor, y hoy debate si debe dar más peso a la escala y a la resiliencia.
\end{itemize}

% --------------------------------------------------------------------------
\subsection{Motivación de la teoría económica}\label{sec:motivacion}

\begin{itemize}
\item Para que el mercado común funcione no basta con quitar barreras (integración negativa): hacen falta políticas comunes que aprovechen su potencial (integración positiva). La de competencia es la principal, junto con otras como la agrícola o la de transportes, que impiden que las políticas nacionales creen ventajas artificiales.
\item El modelo de \ref{sec:modelo-integracion} da la razón: la ganancia del mercado único viene de la competencia. Si las empresas que sobreviven a la racionalización se reparten el mercado, abusan de su tamaño o se fusionan hasta recuperar el margen, la ganancia desaparece; y si los Estados subvencionan a las suyas, la racionalización no se produce.
\item \emph{Los efectos de la competencia}:
  \begin{itemize}
  \item \emph{eficiencia asignativa}: precios cercanos al coste marginal y más variedad \vertema{3.A.17};
  \item \emph{eficiencia productiva}: menos costes y menos ineficiencia X (\autor{Leibenstein}, 1966)\nocite{Leibenstein1966Allocative} \vertema{3.A.12};
  \item \emph{eficiencia dinámica}: más innovación \vertema{3.A.15}.
  \end{itemize}
\item \emph{Las 2~tradiciones}: la de Harvard, que juzgaba por la estructura del mercado, y la de Chicago, que solo veía problema en la conducta que reduce el bienestar del consumidor \vertema{3.A.15}. La Unión añadió la de Friburgo, el ordoliberalismo, para el que la competencia protege también la libertad económica y el propio mercado único. Desde los 2000, con el «enfoque económico más riguroso», la Comisión juzga por los efectos.\footnote{\autor{George Stigler}, de la escuela de Chicago, recibió el Premio Nobel de Economía en 1982 «por sus estudios pioneros de las estructuras industriales, el funcionamiento de los mercados y las causas y efectos de la regulación pública»\nocite{Nobel1982Stigler}, y \autor{Jean Tirole}, en 2014, «por su análisis del poder de mercado y de la regulación»\nocite{Nobel2014Tirole}.}
\end{itemize}

% --------------------------------------------------------------------------
\subsection{Marco jurídico: los 4~pilares}\label{sec:marco-competencia}

\begin{itemize}
\item \emph{Delimitación de competencias}: el establecimiento de las normas sobre competencia necesarias para el funcionamiento del mercado interior es una de las 5~competencias \emph{exclusivas} de la Unión (art. 3.1.b TFUE): solo ella legisla, y los Estados actúan si los faculta o para aplicar sus actos (art. 2.1) (anexo~\ref{anexo:competencias}). Las normas europeas se aplican cuando la conducta puede afectar al comercio entre los Estados; si no, rige el Derecho nacional \vertema{4.A.15}.
\item \emph{La base jurídica} son los arts. 101~a 109~TFUE y los Reglamentos 1/2003 y 139/2004, y se ordena en 4~pilares clásicos a los que se han sumado 2~instrumentos nuevos (figura~\ref{figu:pilares}):
  \begin{itemize}
  \item los \emph{acuerdos} entre empresas (art. 101~y Reglamento 1/2003);
  \item el \emph{abuso de posición dominante} (art. 102);
  \item el control de \emph{concentraciones} (Reglamento 139/2004);
  \item el control de las \emph{ayudas de Estado} (arts. 107-109);
  \item y, desde 2023, las \emph{subvenciones extranjeras} y los \emph{mercados digitales}.
  \end{itemize}
\item \emph{La toma de decisiones}, a 2~niveles:
  \begin{itemize}
  \item \emph{las normas}: el Consejo, a propuesta de la Comisión y previa consulta al Parlamento, adopta por mayoría cualificada los reglamentos que desarrollan los arts. 101~y 102 (art. 103) y las exenciones de ayudas (art. 109); la Comisión adopta los reglamentos de exención por categorías y las directrices;
  \item \emph{los casos}: los decide el Colegio de Comisarios, preparados por la Dirección General de Competencia, bajo la vicepresidenta ejecutiva \autor{Teresa Ribera}; en los arts. 101~y 102~también las autoridades nacionales (en España, la Comisión Nacional de los Mercados y la Competencia, CNMC) y los jueces. Sus decisiones las controlan el Tribunal General y el Tribunal de Justicia \vertema{3.B.36}.
  \end{itemize}
\item \emph{El reparto de los casos} en la Red Europea de Competencia: se presume que un acuerdo no afecta de forma apreciable al comercio entre Estados si la cuota conjunta de las empresas no pasa del 5~\% y su volumen de negocios en la Unión no supera los 40~millones de euros; y un caso lo lleva la autoridad «mejor situada», la que puede ponerle fin y reunir las pruebas donde la conducta tiene efecto. La Comisión lo está cuando los efectos alcanzan a más de 3~Estados\nocite{CE2004Red}.
\item \emph{La evolución}, en 4~etapas:
  \begin{itemize}
  \item \emph{1957-1989, la construcción}: el Tratado de Roma ya prohibía los cárteles, los abusos y las ayudas; el Reglamento 17/1962 obligaba a notificar los acuerdos a la Comisión, que tenía el monopolio de las exenciones;
  \item \emph{1989-2004, el control de concentraciones}: el Reglamento 4064/89, sustituido por el 139/2004 con el test SIEC;
  \item \emph{2004-2019, la modernización}: el Reglamento 1/2003 descentralizó la aplicación y suprimió la notificación, se generalizó la clemencia (desde 1996) y el «enfoque económico más riguroso», y en 2012-2014 se modernizaron las ayudas de Estado;
  \item \emph{desde 2020, los nuevos frentes}: los marcos temporales de ayudas, el Reglamento de Mercados Digitales y el de subvenciones extranjeras (2022) y las directrices sobre abusos (2026) y concentraciones (en preparación).
  \end{itemize}
\end{itemize}

\begin{infografia}[H]
\centering
\caption{Los pilares de la política de competencia de la Unión Europea}\label{figu:pilares}
\grafico[0.98\textwidth]{pilares-competencia}
\fuente{elaboración propia a partir del TFUE y de los reglamentos citados. RGEC: Reglamento general de exención por categorías.}
\end{infografia}

\subsubsection{Primer pilar: los acuerdos entre empresas}\label{sec:acuerdos}

\begin{itemize}
\item \emph{La prohibición} (art. 101.1~TFUE): son incompatibles con el mercado interior y quedan prohibidos los acuerdos entre empresas, las decisiones de asociaciones de empresas y las prácticas concertadas que puedan afectar al comercio entre los Estados y que tengan por objeto o efecto impedir, restringir o falsear la competencia. Por ejemplo:
  \begin{itemize}
  \item fijar los precios u otras condiciones de transacción;
  \item limitar o controlar la producción, los mercados, el desarrollo técnico o las inversiones;
  \item repartirse los mercados o las fuentes de abastecimiento;
  \item discriminar a los clientes o imponer prestaciones suplementarias.
  \end{itemize}
  Las empresas que se coordinan así forman un \concepto{cártel}, y sus acuerdos son nulos de pleno derecho (art. 101.2).
\item \emph{La excepción} (art. 101.3): un acuerdo restrictivo se permite si cumple a la vez 4~condiciones:
  \begin{itemize}
  \item mejora la producción o la distribución o fomenta el progreso técnico o económico;
  \item reserva a los usuarios una participación equitativa en el beneficio;
  \item sus restricciones son indispensables para alcanzar esos objetivos;
  \item no elimina la competencia en una parte sustancial del mercado.
  \end{itemize}
  La aplican los reglamentos de exención por categorías (acuerdos verticales, Reglamento (UE) 2022/720; I+D y especialización, 2023/1066 y 2023/1067) y las directrices horizontales de 2023, con un capítulo sobre los acuerdos de sostenibilidad. Los acuerdos de menor importancia (cuota conjunta de menos del 10~\% entre competidores y del 15~\% entre no competidores) no suelen restringir de forma apreciable la competencia, salvo los cárteles\nocite{CE2014Minimis}.
\item \emph{El procedimiento}: desde el \concepto{Reglamento 1/2003} (aplicable desde mayo de 2004), los acuerdos ya no se notifican: es la «excepción legal», que cada empresa evalúa por sí misma, y los arts. 101~y 102~los aplican también las autoridades nacionales y los jueces, coordinados en la \emph{Red Europea de Competencia}. Las multas llegan al 10~\% del volumen de negocios total de la empresa\nocite{UE2003Reglamento1}. Su revisión, en preparación desde 2024, no tiene aún propuesta.
\item \emph{La clemencia}\nocite{UE2006Clemencia}: la primera empresa que aporta pruebas de un cártel queda exenta de la multa, y las siguientes obtienen reducciones del 30-50~\%, del 20-30~\% y de hasta el 20~\%; con la \emph{transacción}, quien reconoce la infracción obtiene un 10~\% más. La mayor parte de los cárteles se descubren así, porque convierte el cártel en un dilema del prisionero, en el que cada miembro quiere confesar antes que los demás, y exige a sus miembros mucha más paciencia para sostenerlo \eqref{eq:colusion-clemencia} \vertema{3.A.14} \vertema{3.A.19}.
\item \emph{Casos}: los camiones (2016-2017, 3.807~millones de euros, la mayor multa de la historia); los derivados del euríbor (2016, unos 485~millones a Crédit Agricole, HSBC y JPMorgan); y en 2025 el reciclaje de vehículos (458~millones a 15~fabricantes) y \emph{Delivery Hero} y \emph{Glovo} (329~millones), el primer cártel sancionado por pactos para no contratarse los trabajadores\nocite{CE2026Carteles}.
\item \emph{¿Por qué perseguir los cárteles con tanta dureza y cuánto debe ser la multa?}
  \begin{itemize}
  \item \emph{El coste}: un cártel reproduce el monopolio, sube el precio y reduce la cantidad \vertema{3.A.17}. Crea una transferencia de los consumidores a las empresas, el sobreprecio, que es su ganancia ilícita, y una pérdida de eficiencia que no gana nadie, el triángulo de \autor{Arnold Harberger} (1954)\nocite{Harberger1954Monopoly}. Los sobreprecios son grandes: la mediana de los cárteles que estudia \textcite{Connor2014Price} es del 23~\% del precio.
  \item \emph{La multa disuasoria}: siguiendo a \textcite{Becker1968Crime}, que trató el delito como una decisión racional,\footnote{\autor{Gary Becker} recibió el Premio Nobel de Economía en 1992 «por haber extendido el dominio del análisis microeconómico a un amplio abanico de comportamientos e interacciones humanas, incluido el comportamiento fuera del mercado»\nocite{Nobel1992Becker}.} las empresas forman el cártel si lo que ganan supera la multa que esperan pagar \parencite{Landes1983Optimal}. Como solo se descubre una parte de los cárteles, la multa tiene que ser la ganancia ilícita multiplicada por la inversa de la probabilidad de descubrirlos, y debe medirse por esa ganancia, no por la pérdida de eficiencia, que es menor.
  \item \emph{Las 2~palancas}: subir la multa, con el límite de no llevar a la quiebra a la empresa, o subir la probabilidad de descubrir el cártel, que es lo que hace la clemencia.
  \item \emph{La evidencia}: \textcite{CombeMonnierLegal2008Cartels} estiman la probabilidad de detección en torno al 13~\% anual para los cárteles sancionados por la Comisión (una cota superior); con sobreprecios del 20~\%, la multa disuasoria sería varias veces la ganancia ilícita, y el límite del 10~\% del volumen de negocios hace pensar que las multas europeas disuaden poco. Las de la Comisión suman unos 30.600~millones de euros en 162~decisiones entre 1990 y 2025, ajustadas por las sentencias; tras el máximo de 2005-2019, menos de 2.400~millones entre 2020 y 2024 y 860~millones en 2025, en 4~decisiones (gráfico~\ref{fig:multas})\nocite{CE2026Carteles}. Con las directrices de 2006, el importe de base llega al 30~\% de las ventas afectadas por cada año de infracción.
  \item \emph{Lo que se debate}: la caída de las multas desde 2020 puede deberse a que hay menos cárteles o a que se descubren menos, porque la clemencia funciona peor cuando las empresas temen más las demandas de daños de la Directiva 2014/104/UE. La Unión solo multa a las empresas; Estados Unidos y el Reino Unido castigan también con prisión a los directivos.
  \end{itemize}
\end{itemize}

\begin{figure}[H]
\centering
\caption{Multas de la Comisión Europea por cártel, 1990-2025}\label{fig:multas}
\grafico[0.85\textwidth]{multas-carteles}
\fuente{Comisión Europea, DG Competencia, \emph{Cartel statistics} (actualizadas a 26/02/2026)\nocite{CE2026Carteles}. Millones de euros, ajustados por las sentencias de los tribunales, por quinquenios; bajo cada barra, el número de decisiones.}
\end{figure}

\begin{anotaciones}
\textbf{El coste social del cártel y la multa disuasoria.} Con demanda decreciente y coste marginal constante $c$, sin cártel el mercado es competitivo ($p^c=c$, $Q^c$); el cártel actúa como un monopolista, vende $Q^K$, donde $\mathrm{IMg}=c$, y cobra $p^K$ (gráfico~\ref{fig:cartel}):
\begin{itemize}
\item el sobreprecio, $\Delta B=(p^K-c)\,Q^K$, es la ganancia ilícita del cártel;
\item la pérdida irrecuperable de eficiencia es $\mathrm{PIE}=\tfrac12\,(p^K-c)\,(Q^c-Q^K)$.
\end{itemize}
Si el cártel se descubre y se sanciona con probabilidad $\pi$ y paga una multa $M$, a empresas neutrales al riesgo les conviene formarlo si $\Delta B-\pi M>0$. Para disuadirlas, la multa esperada debe superar la ganancia:
\begin{equation}\label{eq:multa-disuasoria}
M\geq\frac{\Delta B}{\pi}
\end{equation}
Si se descubre uno de cada 6~cárteles ($\pi=1/6$), la multa disuasoria es 6~veces la ganancia ilícita.
\begin{figuracaja}
\centering
\caption{El coste social de un cártel: sobreprecio y pérdida de eficiencia}\label{fig:cartel}
\grafico[0.58\textwidth]{cartel-bienestar}
\fuente{elaboración propia a partir de \textcite{Harberger1954Monopoly} y \textcite{Landes1983Optimal}, con $p=8-Q$ y $c=2$: el cártel vende $Q^K=3$ a $p^K=5$; el sobreprecio es 9~y la pérdida de eficiencia, 4,5.}
\end{figuracaja}
\end{anotaciones}

\begin{anotaciones}
\textbf{La clemencia y la sostenibilidad del cártel.} Un cártel no es un equilibrio: a cada miembro le conviene desviarse, y solo se sostiene si el castigo futuro pesa más que la ganancia de hoy. Con $B^{\mathrm C}$ el beneficio por periodo de cada empresa si todas coluden, $B^{\mathrm D}$ el del periodo en que se desvía y $B^{\mathrm N}$ el del equilibrio de \autor{Nash} de cada periodo, que sigue a la desviación como castigo, la colusión es sostenible si las empresas son lo bastante pacientes, es decir, si su factor de descuento $\beta$ cumple \vertema{3.A.19}:
\begin{equation}\label{eq:colusion-clemencia}
\beta\;\geq\;\frac{B^{\mathrm D}-B^{\mathrm C}}{B^{\mathrm D}-B^{\mathrm N}}
\end{equation}
El numerador es lo que se gana hoy desviándose, y el denominador, lo que se pierde cada periodo después. La clemencia sube el umbral de paciencia por 2~vías:
\begin{itemize}
\item \emph{baja $B^{\mathrm C}$}: seguir en el cártel es más arriesgado, porque cualquier socio puede denunciarlo, y la multa esperada se resta del beneficio de coludir;
\item \emph{sube $B^{\mathrm D}$}: quien rompe el acuerdo y lo denuncia primero se queda la ganancia de la desviación y, además, no paga multa.
\end{itemize}
Las 2~vías estrechan la distancia entre lo que da coludir y lo que da el castigo, $B^{\mathrm C}-B^{\mathrm N}$, que es lo que disciplina al cártel, y agrandan la tentación de hoy: hace falta mucha más paciencia para coludir. Con más paciencia necesaria, menos cárteles son sostenibles; y los que existen se vuelven un dilema del prisionero entre sus miembros, que compiten por denunciar los primeros \parencite{MottaPolo2003Leniency} \vertema{3.A.14}.
\end{anotaciones}

\subsubsection{Segundo pilar: el abuso de posición dominante}\label{sec:abuso}

\begin{itemize}
\item \emph{La prohibición} (art. 102~TFUE): es incompatible con el mercado interior y queda prohibida la explotación abusiva, por una o más empresas, de una posición dominante en el mercado interior o en una parte sustancial de él. No se castiga ser dominante, sino abusar: es como el cártel, pero sin necesidad de coordinarse con nadie.
\item \emph{El análisis}, en 3~pasos:
  \begin{itemize}
  \item \emph{definir el mercado relevante}, de producto y geográfico, con el test del monopolista hipotético; lo regula la comunicación de definición de mercado de 8~de febrero de 2024, que añade los mercados digitales y multilaterales \vertema{3.A.20};\footnote{El mercado de producto lo forman los bienes que los consumidores consideran sustituibles por sus características, su precio y su uso; el geográfico, la zona donde las condiciones de competencia son homogéneas. A veces hay que considerar también la dimensión temporal.}
  \item \emph{determinar si hay dominio}: la jurisprudencia lo presume con cuotas de más del 50~\% y lo ve poco probable por debajo del 40~\%;
  \item \emph{determinar si la conducta es abusiva}: \emph{excluyente} (precios predatorios, descuentos de fidelidad, ventas vinculadas, negativa a suministrar, autopreferencia) o \emph{explotadora} (precios excesivos, condiciones no equitativas).
  \end{itemize}
  Las directrices sobre abusos excluyentes, adoptadas el 3~de septiembre de 2026, sustituyen a la comunicación de prioridades de 2009 y dan menos peso al test del «competidor igual de eficiente»\nocite{CE2026Art102}.
\item \emph{Casos}: Google Shopping (2.424~millones de euros, confirmada en 2024), Google Android (4.125~millones, confirmada por el Tribunal de Justicia el 2~de julio de 2026), Apple y la música en \emph{streaming} (1.840~millones, 2024), Meta y Marketplace (798~millones, 2024) y Google y la publicidad en línea (2.950~millones, 2025); en junio de 2026, medidas cautelares contra Meta para que deje a los asistentes de inteligencia artificial rivales acceder a WhatsApp\nocite{TJUE2026GoogleAndroid,CE2026Hechos}.
\item \emph{El Reglamento de Mercados Digitales} (DMA, 2022/1925), aplicable desde 2023, añade una regulación \emph{ex ante}: no espera al abuso, sino que impone obligaciones (interoperabilidad, portabilidad de los datos, prohibición de autopreferencia) a los \concepto{guardianes de acceso}, hoy 7 (Alphabet, Amazon, Apple, Booking, ByteDance, Meta y Microsoft) con 23~servicios básicos de plataforma. Sus primeras multas fueron a Apple (500~millones) y Meta (200~millones) en 2025, y a Google (890~millones) en julio de 2026; la revisión de abril de 2026 concluyó que es pronto para cambiarlo\nocite{UE2022DMA,CE2026DMA} \vertema{3.A.15}.
\end{itemize}

\subsubsection{Tercer pilar: el control de concentraciones}\label{sec:concentraciones}

\begin{itemize}
\item \emph{El principio} (Reglamento 139/2004): se declaran incompatibles las concentraciones que obstaculicen de forma significativa la competencia efectiva en el mercado interior o en una parte sustancial de él, en especial por crear o reforzar una posición dominante (el \emph{test SIEC}) \vertema{3.B.4}\nocite{UE2004Reglamento139}.
\item \emph{La notificación previa} es obligatoria si se superan los umbrales de «dimensión comunitaria»:
  \begin{itemize}
  \item \emph{principales}: un volumen de negocios mundial conjunto de más de 5.000~millones de euros y de más de 250~millones en la Unión para al menos 2~de las empresas;
  \item \emph{alternativos}: 2.500~millones mundiales y más de 100~millones en al menos 3~Estados, con al menos 2~empresas por encima de 25~millones en cada uno;
  \item en ambos casos, salvo que cada empresa haga más de 2~tercios de su negocio en la Unión en un mismo Estado: entonces es asunto nacional.
  \end{itemize}
\item \emph{El procedimiento}:
  \begin{itemize}
  \item \emph{fase I}, de 25~días hábiles: la gran mayoría de las operaciones se autorizan aquí, muchas por el procedimiento simplificado;
  \item \emph{fase II}, de 90~días hábiles, si hay dudas serias: la Comisión autoriza, autoriza con compromisos (vender negocios, dar acceso a los rivales) o prohíbe. Las decisiones se pueden recurrir ante el Tribunal General y, en casación, ante el Tribunal de Justicia.
  \end{itemize}
\item \emph{Las cifras}: 10.200~notificaciones desde 1990, 311~fases II y 33~prohibiciones, la última en 2023 (\emph{Booking}/\emph{eTraveli}); entre ellas, \emph{Siemens}/\emph{Alstom} (2019), que la Comisión prohibió por sus efectos en la señalización ferroviaria y los trenes de muy alta velocidad (gráfico~\ref{fig:concentraciones})\nocite{CE2026Concentraciones}.
\item \emph{Las adquisiciones de empresas pequeñas}: las que no alcanzan los umbrales pueden eliminar a un rival potencial (las \emph{adquisiciones asesinas}). La Comisión las revisaba si se lo pedía un Estado (art. 22), hasta que el Tribunal de Justicia lo limitó en \emph{Illumina}/\emph{Grail} (3~de septiembre de 2024); hoy 9~Estados tienen mecanismos para revisarlas de oficio.
\end{itemize}

\begin{figure}[H]
\centering
\caption{Concentraciones notificadas a la Comisión Europea, 1991-2025}\label{fig:concentraciones}
\grafico[0.85\textwidth]{concentraciones}
\fuente{Comisión Europea, DG Competencia, \emph{Merger statistics} (a 31/10/2026)\nocite{CE2026Concentraciones}. Barras, eje izquierdo: notificaciones por el procedimiento simplificado y ordinarias; línea, eje derecho: fases II abiertas. Desde 1990: 10.200~notificaciones y 33~prohibiciones.}
\end{figure}

\begin{itemize}
\item \emph{¿Cuándo es buena una fusión que da poder de mercado?} Una fusión entre competidores tiene 2~efectos opuestos: reduce la competencia, y el precio puede subir; y puede ahorrar costes por la escala o las sinergias. \textcite{Williamson1968Economies} preguntó cuánto ahorro compensa una subida de precio:\footnote{\autor{Oliver Williamson} recibió el Premio Nobel de Economía en 2009 «por su análisis de la gobernanza económica, especialmente de los límites de la empresa»\nocite{Nobel2009OstromWilliamson}.}
  \begin{itemize}
  \item \emph{la respuesta}: muy poco; como la pérdida de eficiencia depende del \emph{cuadrado} de la subida de precio, y el ahorro de costes se gana sobre toda la producción, con el criterio del \emph{bienestar total} casi cualquier fusión con algún ahorro sería aceptable;
  \item \emph{el matiz de la Unión}: su criterio es el \emph{bienestar del consumidor}, para el que la subida de precio es una transferencia que sí cuenta; por eso solo acepta las eficiencias que sean verificables, propias de la fusión y trasladables a los consumidores, de modo que el precio no suba;
  \item \emph{en oligopolio}: en \autor{Cournot}, una fusión sin sinergias sube siempre el precio, y para que no suba el ahorro de costes marginales tiene que ser grande (\autor{Farrell} y \autor{Shapiro}, 1990)\nocite{FarrellShapiro1990Horizontal} \vertema{3.A.19};
  \item \emph{la innovación}: una fusión puede reducir la competencia por innovar o financiar la innovación que necesita escala; las \emph{adquisiciones asesinas} son el caso extremo: en farmacia, entre el 5~y el 7~\% de las compras sirven para cerrar proyectos que competían con los del comprador\nocite{CunninghamEdererMa2021Killer}.
  \end{itemize}
\item \emph{La evidencia}: las evaluaciones a posteriori de 42~fusiones en Estados Unidos que reúne \textcite{Kwoka2015Mergers} encuentran una subida media de precios del 7~\%, también en las que se autorizaron con la venta de activos: el control fue, en conjunto, demasiado permisivo. En la Unión, las notificaciones rondan las 350-400 al año, con muy pocas fases II (3-4 al año desde 2024) y ninguna prohibición desde 2023 (gráfico~\ref{fig:concentraciones}).
\item \emph{El debate actual}: el criterio de \autor{Williamson} es estático: no recoge la innovación, la resiliencia ni la escala necesaria para competir fuera. De ahí la \emph{revisión de las directrices}, abierta en mayo de 2025: el borrador de 30~de abril de 2026 une las de fusiones horizontales y no horizontales y valora la escala, la resiliencia y la innovación con una «teoría del beneficio»; su adopción está prevista para finales de 2026\nocite{CE2026DirectricesBorrador}. Y el control de concentraciones solo vigila a las empresas: los Estados también pueden cerrar el mercado, ayudando a las suyas.
\end{itemize}

\begin{anotaciones}
\textbf{El compromiso de \autor{Williamson}.} Antes de la fusión, el mercado es competitivo: $p_0=\mathrm{CMe}_0$ y se vende $Q_0$. La fusión baja el coste medio a $\mathrm{CMe}_1$ y da poder de mercado: el precio sube a $p_1$ y se vende $Q_1$ (gráfico~\ref{fig:williamson}). Con $\varepsilon$ la elasticidad de la demanda:
\begin{itemize}
\item la pérdida es el triángulo $A_1\approx\tfrac12\,\Delta p\,\Delta Q$, con $\Delta p=p_1-p_0$ y $\Delta Q=Q_0-Q_1\approx|\varepsilon|\,Q_0\,\Delta p/p_0$;
\item la ganancia es el rectángulo de ahorro de costes, $A_2=\Delta c\,Q_1\approx\Delta c\,Q_0$, con $\Delta c=\mathrm{CMe}_0-\mathrm{CMe}_1$;
\item la fusión mejora el bienestar total si $A_2>A_1$, es decir, con $c_0=\mathrm{CMe}_0=p_0$, si
\begin{equation}\label{eq:williamson}
\frac{\Delta c}{c_0}>\frac{|\varepsilon|}{2}\left(\frac{\Delta p}{p_0}\right)^{2}
\end{equation}
\end{itemize}
Con $|\varepsilon|=1$, una subida del 10~\% se compensa con un ahorro del 0,5~\% del coste; con $|\varepsilon|=2$ y una subida del 20~\%, con un 4~\%. El rectángulo de transferencia, $(p_1-p_0)\,Q_1$, no cuenta con el criterio del bienestar total, pero sí con el del consumidor.
\begin{figuracaja}
\centering
\caption{El compromiso de Williamson: pérdida de eficiencia frente a ahorro de costes}\label{fig:williamson}
\grafico[0.6\textwidth]{williamson}
\fuente{elaboración propia a partir de \textcite{Williamson1968Economies}, con $p=8-Q$, $\mathrm{CMe}_0=p_0=4$, $\mathrm{CMe}_1=3$ y $p_1=5$: $A_1=0{,}5$ y $A_2=3$. La transferencia, 3, no cuenta con el criterio del bienestar total, pero sí con el del consumidor.}
\end{figuracaja}
\end{anotaciones}

\subsubsection{Cuarto pilar: las ayudas de Estado}\label{sec:ayudas}

\begin{itemize}
\item \emph{La prohibición} (art. 107.1~TFUE): son incompatibles con el mercado interior las ayudas otorgadas por los Estados o con fondos estatales, bajo cualquier forma, que falseen o amenacen falsear la competencia, favoreciendo a determinadas empresas o producciones, en la medida en que afecten al comercio entre Estados. Hay ayuda si se cumplen a la vez 4~criterios:
  \begin{itemize}
  \item recursos estatales, imputables al Estado (subvenciones, préstamos o garantías en condiciones mejores que las del mercado, exenciones fiscales, aportaciones de capital);
  \item una ventaja económica para quien la recibe;
  \item selectividad: favorece a unas empresas o sectores y no a todos;
  \item que falsee la competencia y afecte al comercio entre Estados.
  \end{itemize}
\item \emph{Las excepciones}:
  \begin{itemize}
  \item \emph{automáticas} (art. 107.2): ayudas sociales a consumidores individuales, ayudas por catástrofes naturales o acontecimientos excepcionales y las de la división de Alemania;
  \item \emph{posibles} (art. 107.3): el desarrollo de las regiones con un nivel de vida anormalmente bajo o ultraperiféricas, los \concepto{proyectos importantes de interés común europeo} (IPCEI) y la perturbación grave de la economía de un Estado, el desarrollo de actividades o regiones, la cultura y el patrimonio, y las que decida el Consejo.
  \end{itemize}
\item \emph{El procedimiento} (art. 108): los Estados notifican las ayudas y no pueden pagarlas hasta que la Comisión las autoriza (cláusula de suspensión); la ayuda ilegal e incompatible se recupera con intereses. El Consejo puede eximir categorías de ayudas (art. 109).
\item \emph{Las exenciones}: el \emph{Reglamento general de exención por categorías} (RGEC, 651/2014), vigente hasta el 31~de diciembre de 2026 y con uno nuevo previsto para 2027, exime de notificar las ayudas a pymes, I+D, formación, medio ambiente o regionales que cumplen sus condiciones; las de \emph{minimis}, hasta 300.000~euros por empresa en 3~años (Reglamento (UE) 2023/2831), no se consideran ayuda\nocite{UE2023Minimis,CE2026RGEC}.
\item \emph{Los marcos de crisis}, que flexibilizan los criterios de compatibilidad durante un tiempo:
  \begin{itemize}
  \item el marco temporal de la COVID-19 (marzo de 2020);
  \item el de la guerra de Ucrania (marzo de 2022), de crisis y transición desde marzo de 2023;
  \item el \emph{marco del Pacto por una Industria Limpia} (CISAF), desde el 25~de junio de 2025 hasta 2030, para renovables, descarbonización y fabricación de tecnologías limpias.
  \end{itemize}
\item \emph{Los proyectos importantes de interés común europeo}: hay 11~IPCEI aprobados (microelectrónica, baterías, hidrógeno, nube, salud).
\item \emph{Las cifras}: las ayudas pasaron del 0,97~\% del PIB de la Unión en 2019 al 2,4~\% en 2020; en 2024 fueron 168.230~millones de euros, el 0,94~\% del PIB, ya casi todas ordinarias (gráfico~\ref{fig:ayudas}). Alemania concentra una cuarta parte del gasto, y España gastó el 0,99~\% de su PIB\nocite{CE2026Scoreboard}.
\item \emph{Las empresas públicas}: las mismas reglas se aplican a las empresas públicas y a las que prestan servicios de interés económico general, con requisitos de transparencia de sus relaciones financieras con el Estado (art. 106~TFUE) \vertema{4.B.15}.
\end{itemize}

\begin{figure}[H]
\centering
\caption{Ayudas de Estado en la Unión Europea, 2000-2024}\label{fig:ayudas}
\grafico[0.98\textwidth]{ayudas-estado}
\fuente{Comisión Europea, DG Competencia, \emph{State aid Scoreboard 2025} (datos de 2024; consulta del 14/11/2026)\nocite{CE2026Scoreboard}. Elemento de ayuda en porcentaje del PIB, sin ferrocarriles ni ayudas al sector financiero. Crisis: marcos temporales de la COVID-19 y de la guerra de Ucrania.}
\end{figure}

\begin{itemize}
\item \emph{¿Por qué tiene que controlar la Unión las ayudas, si cada Estado gasta su propio dinero?} Porque, en un mercado integrado, la ayuda de un Estado desplaza la producción y el empleo de los demás:
  \begin{itemize}
  \item \emph{la carrera de subvenciones}: a cada Estado le conviene ayudar a su empresa, haga lo que haga el otro, porque le quita cuota a la rival; si todos lo hacen, ninguno gana cuota y todos pagan la ayuda con impuestos que distorsionan. Es un \emph{dilema del prisionero} entre Estados, con la misma estructura que el de la política comercial estratégica \vertema{3.A.14} \vertema{3.B.8};
  \item \emph{la solución}: la prohibición de las ayudas es la solución cooperativa que los Estados no alcanzarían solos, y la Comisión, como autoridad común, la hace cumplir; \textcite{Collie2000State} muestra que, en un mercado integrado y con impuestos que distorsionan, prohibirlas mejora el bienestar del conjunto;
  \item \emph{la asimetría fiscal}: si unos Estados pueden gastar más que otros, el juego deja de ser simétrico, y las empresas de los Estados con más margen fiscal ganan cuota aunque no sean las más eficientes; es la crítica a los marcos temporales, cuya ayuda aprobada se concentró en buena parte en Alemania y Francia;
  \item \emph{el límite del argumento}: justifica que el control sea europeo, pero no dice cuánta ayuda es buena; las ayudas sí se justifican cuando corrigen un fallo de mercado (I+D, medio ambiente, regiones)\nocite{BesleySeabright1999Effects} \vertema{3.A.23}.
  \end{itemize}
\item \emph{La evidencia}: la ayuda de crisis llegó al 1,4~\% del PIB de la Unión en 2020, y la ayuda total en 2024 va del 1,37~\% del PIB de Hungría al 0,40~\% de Portugal e Irlanda (gráfico~\ref{fig:ayudas}). Y el juego ya no es solo entre Estados de la Unión: las empresas de terceros países también reciben subvenciones de sus gobiernos.
\end{itemize}

\begin{anotaciones}
\textbf{La carrera de subvenciones como dilema del prisionero.} 2~Estados, cada uno con una empresa que compite con la del otro en el mercado único, deciden si la subvencionan; el bienestar de cada uno es el beneficio de su empresa menos el coste de la ayuda. Con unos pagos ilustrativos (tabla~\ref{tab:subvenciones}):
\begin{itemize}
\item si el otro no subvenciona, a cada Estado le conviene subvencionar (14~frente a 10);
\item si el otro subvenciona, también (7~frente a 5);
\item subvencionar es una estrategia dominante, y el equilibrio de \autor{Nash} es que los 2~subvencionen (7, 7), peor para los 2~que no hacerlo (10, 10).
\end{itemize}
\begin{tablacaja}
\centering
\caption{La carrera de subvenciones como dilema del prisionero (bienestar de cada Estado)}\label{tab:subvenciones}
\begin{tabular}{l|c|c|}
\multicolumn{1}{c}{} & \multicolumn{1}{c}{Estado 2 no subvenciona} & \multicolumn{1}{c}{Estado 2 subvenciona}\\
\hhline{~|-|-|}
Estado 1 no subvenciona & \textcolor{tceeverde}{\textbf{(10, 10)}} & (5, 14)\\
\hhline{~|-|-|}
Estado 1 subvenciona & (14, 5) & \textcolor{tceerojo}{\textbf{(7, 7)}}\\
\hhline{~|-|-|}
\end{tabular}
\fuente{elaboración propia, con pagos ilustrativos. En rojo, el equilibrio de Nash; en verde, el resultado cooperativo que busca el control de ayudas.}
\end{tablacaja}
\end{anotaciones}

\paragraph{Los instrumentos nuevos: las subvenciones extranjeras}

\begin{itemize}
\item \emph{El Reglamento (UE) 2022/2560 sobre subvenciones extranjeras}, aplicable desde el 12~de julio de 2023, extiende la lógica de las ayudas de Estado a las que dan los gobiernos de terceros países\nocite{UE2022Subvenciones}:
  \begin{itemize}
  \item obliga a notificar las adquisiciones de empresas con al menos 500~millones de euros de negocio en la Unión, si el comprador recibió más de 50~millones de contribuciones extranjeras en 3~años, y las ofertas a contratos públicos de al menos 250~millones;
  \item la Comisión puede investigar de oficio y exigir compromisos o prohibir la operación;
  \item casos: el tren búlgaro de CRRC (2024, oferta retirada), \emph{e\&}/\emph{PPF Telecom} (2024), \emph{ADNOC}/\emph{Covestro} (2025) y \emph{JD.com}/\emph{Ceconomy} (noviembre de 2026), autorizadas con compromisos; 273~notificaciones hasta mayo de 2026\nocite{CE2026Subvenciones}.
  \end{itemize}
\end{itemize}

% --------------------------------------------------------------------------
\subsection{Promoción, coordinación y debate actual}\label{sec:promocion}

\begin{itemize}
\item \emph{La promoción de la competencia}: la política no solo reprime; también abre mercados, sobre todo las \concepto{industrias de red} en monopolio público:
  \begin{itemize}
  \item \emph{telecomunicaciones}: liberalizadas desde 1998, con el fin del \emph{roaming} en 2017;
  \item \emph{energía}: 4~paquetes (1996-2019) separan la red de la generación y la comercialización y dan acceso a terceros, con agencias reguladoras independientes \vertema{4.A.5} \vertema{4.A.6};
  \item \emph{transporte}: el aéreo, liberalizado en los 90 (cielo único europeo); por carretera, con límites al cabotaje; el ferroviario, con 4~paquetes, el último abierto a la competencia de viajeros desde 2020 \vertema{4.A.12}.
  \end{itemize}
  En todas, el método es separar la red, que es un monopolio natural, de los servicios, que pueden competir \vertema{3.A.20}.
\item \emph{La coordinación}:
  \begin{itemize}
  \item \emph{dentro de la Unión}, la Red Europea de Competencia reparte los casos entre la Comisión y las autoridades nacionales, cuya independencia y poderes reforzó la Directiva (UE) 2019/1 (ECN+);
  \item \emph{fuera}, acuerdos de cooperación con Estados Unidos, Canadá, Japón y otros, y la Red Internacional de Competencia, de la que la Unión es fundadora. No hay reglas de competencia mundiales: la OMC solo tiene las de subvenciones \vertema{3.B.33}.\footnote{Si el mercado relevante es cada vez más mundial, harían falta normas de competencia también mundiales. La OMC solo lo ha logrado en parte con el Acuerdo sobre Subvenciones y Medidas Compensatorias (con poco éxito en agricultura y en las notificaciones) y, de forma plurilateral, en la contratación pública.}
  \end{itemize}
\item \emph{El debate de los «campeones europeos»}: tras la prohibición de \emph{Siemens}/\emph{Alstom} (2019), Francia y Alemania pidieron reglas que permitieran crear grandes empresas europeas capaces de competir con las de China y Estados Unidos, y el informe \autor{Draghi} (2024) pidió valorar la innovación y la escala\nocite{Draghi2024Future}:
  \begin{itemize}
  \item \emph{a favor}: la escala importa en sectores de costes fijos enormes (semiconductores, defensa, telecomunicaciones), y el mercado relevante es cada vez más mundial;
  \item \emph{en contra}: el campeón protegido sube los precios en casa y no compite mejor fuera; la evidencia muestra que los mercados europeos se volvieron más competitivos que los estadounidenses gracias a su política de competencia \nocite{GutierrezPhilippon2023How}, y lo que impide crecer a las empresas europeas es la fragmentación del mercado, no el control de fusiones.
  \end{itemize}
\item \emph{Recapitulación del bloque~\ref{sec:competencia}}: la política de competencia protege la ganancia del mercado único frente a empresas que se coordinan, abusan o se fusionan y frente a Estados que subvencionan. La multa disuasoria debe ser la ganancia entre la probabilidad de detección, una fusión se justifica si el ahorro supera la pérdida y el control de ayudas resuelve un dilema del prisionero entre Estados.
\end{itemize}

% ==========================================================================
\section*{Conclusión}
\minutos{1,5}

\begin{itemize}
\item \textbf{\emph{Recapitulación}}: En conclusión,
  \begin{itemize}
  \item El mercado único es el núcleo económico de la Unión: se construyó en 2~tiempos, la unión aduanera de 1968 y el programa de 1992, y su ganancia principal es la de la competencia y la escala: con mercados más grandes, márgenes más pequeños y empresas más eficientes.
  \item Sus 4~libertades han avanzado de forma desigual: mucho en bienes y capitales, poco en servicios y personas, donde el FMI estima barreras equivalentes a aranceles muy altos.
  \item La política de competencia protege esa ganancia frente a las empresas (cárteles, abusos, concentraciones) y frente a los Estados (ayudas), y desde 2023 frente a las subvenciones extranjeras y los guardianes digitales.
  \end{itemize}
\item \textbf{\emph{Relevancia}}:
  \begin{itemize}
  \item Para España, el mercado único es su principal mercado exterior y el marco de su política de competencia, de ayudas y de buena parte de su regulación \vertema{4.A.15} \vertema{4.A.22}.
  \end{itemize}
\item \textbf{\emph{Extensiones y relación con otras partes del temario}}:
  \begin{itemize}
  \item La teoría de la integración \vertema{3.B.10}; las instituciones que hacen las reglas \vertema{3.B.36}; la política comercial frente a terceros \vertema{3.B.41}; la cohesión, que compensa a la periferia \vertema{3.B.40}; la unión bancaria y la de mercados de capitales \vertema{3.B.44}; y la defensa de la competencia en España \vertema{4.A.15}.
  \end{itemize}
\item \textbf{\emph{Opinión}}:
  \begin{itemize}
  \item La prioridad no es relajar el control de fusiones para fabricar campeones, sino completar el mercado de servicios, de capitales y digital, para que las empresas europeas crezcan compitiendo.

  \item Las directrices nuevas pueden dar más peso a la innovación y a la resiliencia sin renunciar al bienestar del consumidor como criterio.

  \item Quien pide campeones suele ser el que aspira a serlo, y ya lo advirtió \autor{George Stigler} (1971): «por regla general, la industria adquiere la regulación, que se diseña y se aplica sobre todo en su beneficio»\footnote{En el original, \textcite{Stigler1971Theory}, p.~3: «as a rule, regulation is acquired by the industry and is designed and operated primarily for its benefit».}.
  \end{itemize}
\item \textbf{\emph{Idea final}}:
  \begin{itemize}
  \item Volvamos al refrán: el pez grande se come al chico. El mercado único agranda el estanque para que los peces crezcan, y la política de competencia vigila que crezcan por ser mejores, no porque se coman a los demás o porque los alimente su Estado. Si los peces europeos siguen siendo pequeños, la respuesta no está en dejarles comer a los chicos, sino en terminar de quitar las redes que siguen partiendo el estanque.
  \end{itemize}
\end{itemize}

\clearpage
\appendix

% ==========================================================================
\section{Anexo: el modelo de la integración\cortetitulo{} de mercados oligopolísticos}\label{anexo:integracion}

\paragraph{a) Idea}

\begin{itemize}
\item Con mercados nacionales separados, cada uno es pequeño: caben pocas empresas, cada una produce poco (lejos de su escala eficiente) y, como compite con pocas rivales, cobra un margen alto.
\item Al integrarse los mercados, cada empresa compite con más rivales y percibe una demanda más elástica: baja el margen (\concepto{efecto procompetitivo}). Con márgenes más bajos no caben todas las empresas: salen algunas y las que quedan producen más y más barato (\concepto{efecto de racionalización}).
\item Es la lógica de \textcite{SmithVenables1988Completing}, que simularon el mercado único con oligopolios de \autor{Cournot}, y de \textcite{Krugman1979Increasing} con competencia monopolística. La versión más sencilla, con libre entrada, cabe en una fórmula.
\end{itemize}

\paragraph{b) Modelo}

\subparagraph{i) Supuestos}

\begin{enumerate}
\item \emph{Un bien homogéneo} y mercados de tamaño $S$, con demanda inversa lineal $p(Q)=a-bQ$ y $b=1/S$: un mercado el doble de grande tiene, a cada precio, el doble de compradores.
\item \emph{Empresas iguales} con coste $C(q)=\mathrm{CF}+cq$: un coste fijo $\mathrm{CF}$ (economías de escala) y un coste marginal $c$ constante.
\item \emph{Competencia à la \autor{Cournot}} con $n$ empresas \vertema{3.A.19}.
\item \emph{Libre entrada}: entran empresas mientras el beneficio sea positivo (se trata $n$ como continuo).
\item \emph{2 regímenes}: \emph{segmentado}, con 2~países de tamaño $S$ separados por barreras prohibitivas; e \emph{integrado}, un solo mercado de tamaño $2S$ sin barreras.
\end{enumerate}

\subparagraph{ii) Desarrollo}

\begin{itemize}
\item \emph{Paso 1: el equilibrio de \autor{Cournot}} con $n$ empresas:
  \begin{itemize}
  \item cada empresa produce $q_j=\dfrac{a-c}{(n+1)\,b}$ y el precio es $p=\dfrac{a+nc}{n+1}$;
  \item el margen es $p-c=\dfrac{a-c}{n+1}=b\,q_j$: más empresas, menos margen;
  \item el índice de \autor{Lerner} es $\mathcal{L}=\dfrac{p-c}{p}=\dfrac{1}{n\,|\varepsilon|}$, con $\varepsilon$ la elasticidad de la demanda del mercado \vertema{3.A.17}.
  \end{itemize}
\item \emph{Paso 2: la libre entrada}. El beneficio de cada empresa es $(p-c)\,q_j-\mathrm{CF}=b\,q_j^2-\mathrm{CF}$, y la entrada lo lleva a cero:
  \begin{itemize}
  \item $q^{*}=\sqrt{\mathrm{CF}\cdot S}$: cada empresa produce más cuanto mayor es el mercado;
  \item $n^{*}+1=(a-c)\sqrt{S/\mathrm{CF}}$: el número de empresas crece con la raíz del tamaño, no en proporción;
  \item en el equilibrio, el precio iguala el coste medio, $p^{*}=\mathrm{CMe}(q^{*})=c+\mathrm{CF}/q^{*}$, y la demanda residual de cada empresa, de pendiente $-b$, es tangente a la curva de coste medio (gráfico~\ref{fig:mercado-integrado}).
  \end{itemize}
\item \emph{Paso 3: el margen de equilibrio}. Sustituyendo $q^{*}$ en $p-c=b\,q^{*}$:
  \[
  p^{*}-c=\sqrt{\frac{\mathrm{CF}}{S}},
  \]
  que es el resultado \eqref{eq:margen-tamano}: el margen sobre el coste marginal baja con la raíz cuadrada del tamaño del mercado.
\item \emph{Paso 4: de 2~mercados segmentados a uno integrado} ($S\to2S$):
  \begin{itemize}
  \item el margen cae a $1/\sqrt2$ de su valor, un 29~\% menos (efecto procompetitivo);
  \item cada empresa produce $\sqrt2$ veces más, un 41~\% más, y su coste medio baja (racionalización);
  \item en el mercado integrado hay $\sqrt2\,(n^{*}+1)-1$ empresas, menos que las $2n^{*}$ que había sumando los 2~países: unas desaparecen o se fusionan, pero cada consumidor puede comprar a más empresas que antes.
  \end{itemize}
\end{itemize}

\begin{figure}[H]
\centering
\caption{Efecto procompetitivo y de racionalización: el equilibrio de cada empresa en el mercado nacional y en el mercado único}\label{fig:mercado-integrado}
\grafico[0.72\textwidth]{mercado-integrado}
\fuente{elaboración propia a partir de \textcite{SmithVenables1988Completing}, con $c=1$ y $\mathrm{CF}=4$: en el mercado nacional ($S=1$), $q_N=2$ y $p_N=3$; en el mercado único ($S=2$), $q_U=2{,}83$ y $p_U=2{,}41$. $d_N$ y $d_U$ son las demandas residuales de cada empresa, tangentes a la curva de coste medio; la del mercado único es más plana.}
\end{figure}

\subparagraph{iii) Implicaciones}

\begin{itemize}
\item \emph{La ganancia del mercado único es de eficiencia productiva y asignativa a la vez}: precios más bajos para los consumidores y costes medios más bajos por la escala. Por eso \autor{Cecchini} daba ganancias varias veces mayores que las de la unión aduanera.
\item \emph{La integración trae concentración}: menos empresas y más grandes. No es un fallo, sino el mecanismo por el que se aprovecha la escala; pero exige una política de competencia que impida que las empresas que quedan se pongan de acuerdo o se fusionen hasta recuperar el margen (bloque~\ref{sec:competencia}).
\item \emph{Las barreras que no son aranceles también segmentan}: normas técnicas distintas, licencias nacionales o compras públicas cerradas mantienen el $S$ de cada empresa en el tamaño de su país. Quitarlas tiene el mismo efecto que el paso 4.
\end{itemize}

\paragraph{c) Extensiones}

\begin{itemize}
\item \emph{Segmentación frente a integración}: \textcite{SmithVenables1988Completing} distinguen la reducción de los costes de comerciar, que acerca los mercados, de la integración plena, que impide a las empresas fijar precios distintos en cada país. Con mercados todavía segmentados, el comercio cruzado de un mismo bien (el «\emph{dumping} recíproco» de \autor{Brander} y \autor{Krugman}, 1983) ya reduce márgenes; pero las mayores ganancias de sus simulaciones aparecen cuando las empresas dejan de tratar cada país como un mercado aparte \vertema{3.B.10}.
\item \emph{Variedad}: con productos diferenciados, el mercado más grande ofrece además más variedades a cada consumidor \vertema{3.B.6}.
\item \emph{Empresas heterogéneas}: la integración expulsa a las empresas menos productivas y reasigna la producción hacia las mejores, lo que eleva la productividad media \vertema{3.B.6}.
\item \emph{Ineficiencia X}: con más competencia, las empresas también reducen la ineficiencia interna de su organización (\autor{Leibenstein}, 1966) \vertema{3.A.12}.
\end{itemize}

\paragraph{d) Evidencia empírica}

\begin{itemize}
\item \emph{Comercio}: el comercio de bienes dentro de la Unión pasó del 17~\% del PIB en 2002 al 22~\% en 2025, con un pico del 26~\% en 2022 por los precios de la energía; el de servicios, del 4,8~al 8~\% entre 2010 y 2025 (gráfico~\ref{fig:comercio-intra}). En el último informe anual del mercado único, el comercio intra-UE en porcentaje del PIB es uno de los indicadores que empeoran\nocite{CE2026InformeAnual}.
\item \emph{Márgenes}: \textcite{Badinger2007Has} estima que los márgenes sobre el coste marginal bajaron tras el programa de 1992 en la industria y la construcción, pero no en buena parte de los servicios, donde el mercado siguió fragmentado.
\item \emph{PIB}: las estimaciones a posteriori van desde el 3-4~\% del PIB de la Comisión\nocite{CE2025EstrategiaFicha} hasta en torno al 9~\% de los modelos que suman bienes y servicios\nocite{inTVeld2019Economic}; \textcite{MayerVicardZignago2019Cost} estiman que el mercado único ha más que duplicado el comercio de bienes entre sus miembros y que reducirlo a un acuerdo comercial convencional costaría a cada Estado, de media, un 4,4~\% de bienestar (entre el 3~y el 7~\% según el escenario).
\item \emph{Lo que queda}: el FMI estima que en 2020 las barreras dentro de Europa equivalían a un arancel del 44~\% en la manufactura media (frente al 15~\% entre los estados de Estados Unidos) y de hasta el 110~\% en los servicios, y que llevarlas al nivel estadounidense elevaría la productividad casi un 7~\% a largo plazo\nocite{FMI2024REO,FMI2024Kammer}.
\end{itemize}


\paragraph{e) Valoración}

\begin{itemize}
\item El modelo explica bien por qué la ganancia del mercado único se mide en competencia y escala y por qué la integración trae concentración.
\item Pero supone que el mercado integrado sigue siendo competitivo: no dice nada de cárteles, de abusos de las empresas dominantes, de fusiones que eliminan rivales ni de Estados que subvencionan a sus empresas para que sobrevivan a la racionalización. Esas son, precisamente, las tareas de la política de competencia (bloque~\ref{sec:competencia}).
\item Y trata el mercado como un todo, cuando la integración ha avanzado de forma muy desigual según lo que circula: es lo que veremos libertad a libertad.
\end{itemize}


% ==========================================================================
\section{Anexo: el reparto de competencias\cortetitulo{} de la Unión}\label{anexo:competencias}

\begin{itemize}
\item \emph{La regla de partida es la atribución}: la Unión solo tiene las competencias que le dan los Tratados, y las demás siguen siendo de los Estados (arts.~4.1~y 5~TUE). El art.~2~TFUE las ordena en categorías según quién legisla, y los arts.~3~a 6~enumeran los ámbitos de cada una (figura~\ref{figu:competencias-ue}). Las ejerce con \emph{proporcionalidad} (no más de lo necesario) y, fuera de las exclusivas, con \emph{subsidiariedad} (solo si el objetivo se alcanza mejor en común).
\item \emph{Exclusivas y compartidas}: en las \concepto{exclusivas} (art.~3) solo legisla la Unión; los Estados, solo si ella los faculta o para aplicar sus actos. En las \concepto{compartidas} (art.~4), la categoría residual, legislan las 2, pero lo que regula un acto de la Unión deja de estar disponible para los Estados: es la \emph{preferencia} u ocupación del campo (art.~2.2~y Protocolo n.º~25), que solo alcanza a lo que el acto regula. En investigación, espacio, cooperación para el desarrollo y ayuda humanitaria son \emph{paralelas}: la acción de la Unión no impide la de los Estados (art.~4.3~y 4.4).
\item \emph{Coordinación y apoyo}: en las políticas económicas y de empleo (art.~5) la política sigue siendo de los Estados y la Unión la coordina con orientaciones del Consejo; en los ámbitos del art.~6 (salud, industria, cultura, turismo, educación, protección civil, cooperación administrativa) la Unión solo apoya, coordina o complementa, sin armonizar las leyes nacionales. La PESC tiene un régimen propio, intergubernamental, con unanimidad y sin actos legislativos (art.~24~TUE).
\item \emph{Dónde está este tema}: el mercado interior es compartido (art.~4.2.a; arts.~26~y 114~TFUE); las normas sobre competencia necesarias para su funcionamiento (arts.~101~a 109) y la unión aduanera (arts.~28~a 32), en las que se apoya la libre circulación de mercancías, son exclusivas (art.~3.1.a y b).
\item \emph{Qué se deduce}: la exclusiva es legislar las normas de competencia, no aplicarlas: desde el Reglamento 1/2003 también las autoridades y los jueces nacionales aplican los arts.~101~y 102. En el mercado interior, cada directiva o reglamento de armonización ocupa su campo, y los Estados conservan lo que no regula.
\end{itemize}

\begin{infografia}[H]
\centering
\caption{El reparto de competencias de la Unión (arts.~2~a 6~TFUE): el mercado interior y la competencia}\label{figu:competencias-ue}
\grafico[\textwidth]{competencias-ue}
\fuente{elaboración propia a partir de los arts.~2~a 6~del TFUE y de los arts.~4, 5~y 24~del TUE\nocite{UE2016TFUE}.}
\end{infografia}

% ==========================================================================
\section{Anexo: preguntas de test}\label{anexo:test}

\begin{anotaciones}
Preguntas tipo test del primer ejercicio: \href{https://www.victorgutierrezmarcos.es/oposicion/temario/primer-ejercicio/test/simulador.html?tema=3.B.39}{simulador de test}.
\end{anotaciones}

\bibliografia
\end{document}
